% =============================================================================
% CATALOGUE MERGING STRATEGIES
% Self-contained section; include in main.tex with % =============================================================================
% CATALOGUE MERGING STRATEGIES
% Self-contained section; include in main.tex with % =============================================================================
% CATALOGUE MERGING STRATEGIES
% Self-contained section; include in main.tex with % =============================================================================
% CATALOGUE MERGING STRATEGIES
% Self-contained section; include in main.tex with \input{merge_strategies}
% Uses natbib \citep / \citet, booktabs, amsmath, and cleveref already loaded
% in main.tex.
% =============================================================================

\section{Strategies for Merging Earthquake \catalogs}
\label{sec:merging}

Combining records from multiple seismic networks into a single unified \catalog\
is a prerequisite for seismic-hazard analysis, time-series studies of regional
seismicity, and magnitude-frequency assessments that draw on heterogeneous data
archives \citep{woessner2010instrumental, corssa_woessner}.  The task
decomposes into two conceptually distinct steps that must be treated
separately in any robust implementation.

\begin{enumerate}
  \item \textbf{Event association and deduplication.}  Determining which records
        across source \catalogs\ describe the \emph{same physical earthquake}
        and grouping them into clusters.

  \item \textbf{Parameter selection and homogenisation.}  From each cluster,
        selecting or computing a single best origin time, hypocenter, and
        magnitude, and converting all magnitude estimates to a common scale.
\end{enumerate}

Conflating these steps—for instance by applying a magnitude conversion formula
\emph{before} deciding whether two records are duplicates—introduces systematic
biases that propagate into downstream analyses
\citep{Storchak2013, Bondár2011}.

% -----------------------------------------------------------------------------
\subsection{Event Association: Spatio-Temporal Matching}
\label{sec:matching}

\subsubsection{Window-based forward-sweep (classical approach)}

The oldest and most widely deployed method sorts all incoming origin reports by
origin time and scans forward with a sliding window.  If two reports from
\emph{different} networks fall within a fixed time window $\Delta t_{\max}$ and
within a fixed epicentral distance $\Delta d_{\max}$, they are flagged as
candidate duplicates \citep{NCEDCComCat}.

\paragraph{Operational thresholds.}
\Cref{tab:windows} summarises documented threshold choices from major
operational systems.

\begin{table}[H]
  \centering
  \caption{Spatio-temporal matching thresholds documented in major operational
           catalogue systems.  $\Delta t$ = time window, $\Delta d$ = distance
           window.  ANSS thresholds apply to inter-network pairs only;
           intra-network deduplication is assumed upstream.}
  \label{tab:windows}
  \begin{tabular}{llll}
    \toprule
    System & $\Delta t$ (s) & $\Delta d$ (km) & Reference \\
    \midrule
    ANSS Composite \catalog & 16 & 100 & \citet{NCEDCComCat} \\
    Typical regional build ($M \geq 4$) & 60 & 100 & \citet{Storchak2013} \\
    Local / small-event ($M < 4$) & 30 & 30 & \citet{Woessner2005} \\
    Pre-digital era (pre-1960) & 300 & 200 & \citet{Bondár2011} \\
    Dead Sea Transform region & 60 & 100 & \citet{DeadSea2022} \\
    Mexico unified \catalog & 60 & $\approx$111 ($1^{\circ}$) & \citet{Zuniga2019} \\
    \bottomrule
  \end{tabular}
\end{table}

The 16\,s / 100\,km rule of the ANSS Composite \catalog\ is the most precisely
documented public threshold for an active operational system
\citep{NCEDCComCat}.  For global teleseismic events ($M \geq 4.5$) the broader
60\,s / 100\,km window is standard; the ISC bulletin uses a comparable
proximity criterion internally when auto-grouping agency solutions prior to
relocation \citep{ISCBulletinRebuild2020}.

\paragraph{Adaptive thresholds.}
Fixed windows perform poorly across the full magnitude range because larger
events are subject to greater timing and location uncertainty at distant
networks.  The standard remedy, used in the ISC-GEM pipeline and adopted in the
current implementation, is to scale the effective window by the event magnitude
and depth \citep{ISCGEM2018}:

\begin{equation}
  \Delta t_{\mathrm{eff}} = \Delta t_{\mathrm{config}} \cdot f_t(M),
  \qquad
  \Delta d_{\mathrm{eff}} = \Delta d_{\mathrm{config}} \cdot f_d(M) \cdot g(h),
  \label{eq:adaptive}
\end{equation}

where $f_t$, $f_d$ are magnitude multipliers (e.g., 1.0 for $M < 4$, 1.5 for
$4 \leq M < 5.5$, 2.5 for $5.5 \leq M < 7$, 4.0 for $M \geq 7$) and $g(h)$
is a depth multiplier (1.0 for $h < 100$\,km, 1.2 for $100 \leq h < 300$\,km,
1.5 for $h \geq 300$\,km), consistent with ISC-GEM practice
\citep{ISCGEM2018, Engdahl2020}.

\paragraph{Limitation: dense aftershock sequences.}
During intense aftershock sequences the time gap between consecutive legitimate
aftershocks can fall well below 60\,s.  \citet{Tanaka2022} demonstrated that in
the 24 hours following the 2011 Tōhoku $M_\mathrm{w}$\,9.0 earthquake, more
than 599 genuine aftershock pairs had inter-event times under 60\,s—the same
interval used as a duplicate window.  Fixed-window methods therefore cannot
distinguish a duplicate pair from an aftershock pair in such sequences without
additional information.

% - - - - - - - - - - - - - - - - - - - - - - - - - - - - - - - - - - - - - -
\subsubsection{Nearest-neighbour proximity function}
\label{sec:nn}

\citet{Tanaka2022} introduced a physics-based proximity function that
discriminates duplicates from aftershocks by exploiting their different
topological signatures in event-pair space: duplicates form \emph{pairs}
(reports of one earthquake), whereas aftershocks form \emph{branching trees}
(one parent, many offspring).

The proximity between two events $i$ and $j$ is defined as

\begin{equation}
  \eta(i,j) = t_{ij}\; r_{ij}^{d_f}\; 10^{-b m_i},
  \label{eq:proximity}
\end{equation}

where $t_{ij}$ is the inter-event time, $r_{ij}$ is the epicentral distance,
$d_f \approx 1.6$ is the fractal dimension of the epicentre distribution,
$b \approx 1.0$ is the Gutenberg-Richter $b$-value, and $m_i$ is the magnitude
of the preceding event.  Events whose nearest-neighbour distance in this space
falls \emph{below} a threshold derived from network location error (rather than
from earthquake-physics considerations) are classified as duplicates.  Applied
to the Tōhoku sequence, duplicate identification reliability exceeded 97\,\%
\citep{Tanaka2022}.

The method requires estimates of $d_f$ and $b$ for the target seismicity, both
of which can be obtained from the source \catalogs\ themselves via standard
routines \citep{Aki1965, Wiemer2000}.

% - - - - - - - - - - - - - - - - - - - - - - - - - - - - - - - - - - - - - -
\subsubsection{Graph-based connected-components association}
\label{sec:graph}

The methodologically strongest approach constructs an undirected association
graph in which each origin report is a node and each candidate same-event pair
is an edge \citep{Benz2019}.  Connected components of this graph—found in
$O(n\,\alpha(n))$ time via the Union-Find algorithm, where $\alpha$ is the
inverse Ackermann function—correspond to clusters of reports for the same
physical earthquake.

\paragraph{Algorithm sketch.}
\begin{enumerate}
  \item Sort all reports by origin time; build a spatial index (R-tree or
        uniform grid).
  \item For each report $E$, query the index for all reports within
        $\Delta t_{\max}$ and $\Delta d_{\max}$ from \emph{different} source
        networks.
  \item Add a graph edge for each qualifying pair $(E, F)$.
  \item Find connected components; each component is one physical earthquake.
  \item Apply a preference-selection step within each component
        (\cref{sec:selection}).
\end{enumerate}

Compared with the greedy forward sweep, the graph approach naturally handles
\emph{transitive associations}: if network A matches network B, and B matches
network C, all three are placed in the same component even if A and C do not
independently satisfy the proximity criterion.  The greedy sweep can miss such
cases when B has already been marked as processed before A is examined.

\paragraph{Integer linear programming refinement.}
\citet{Benz2019} further applied Integer Linear Programming (ILP) to the
connected-component solution to resolve cases where a component might span two
closely spaced genuine earthquakes.  The ILP minimises the number of source
events required to explain all associations subject to consistency constraints.
Applied to Chile 2010--2017, the method detected 817\,000+ earthquakes at
greater than 95\,\% accuracy on synthetic validation data.

The USGS NEIC Hydra production system uses a graph-association framework
conceptually similar to \citet{Benz2019} \citep{usgs_comcat}.

% -----------------------------------------------------------------------------
\subsection{The ISC-GEM Reference Methodology}
\label{sec:iscgem}

The ISC-GEM Global Instrumental Earthquake \catalog\ \citep{Storchak2013,
ISCGEM2018} covers 1904--2021 for $M \geq 5.5$ globally and $M \geq 5.0$ over
continental regions.  It represents the most thoroughly documented end-to-end
merging pipeline in the seismological literature and serves as the benchmark
against which other implementations are measured.

\subsubsection{Event association}
The ISC automated grouping program compares origin times and hypocentre
locations from all contributing agencies and assembles groups using a
proximity-scoring system that weights azimuthal gap, distance range, station
count, and phase count \citep{Bondár2011}.  Groups may be split or new events
created where the evidence supports distinct earthquakes.

\subsubsection{Relocation}
The ISC applies a two-tier relocation to all grouped events:
\begin{enumerate}
  \item \textbf{EHB algorithm} \citep{Engdahl1998}: improved hypocenters with
        particular attention to focal depth, using reidentified phases.
  \item \textbf{ISC location algorithm} \citep{Bondár2011}: fixes the depth
        output from EHB and reduces systematic location bias by modelling
        correlated travel-time prediction errors across the station network.
\end{enumerate}
The ISC-authored solution is designated \emph{prime} when sufficient phase data
are available.  All agency solutions are retained as secondary hypocenters,
preserving full audit traceability \citep{ISCBulletinRebuild2020}.

\subsubsection{Magnitude assembly}
ISC-GEM reports a single $M_\mathrm{w}$ per event, derived by \citep{ISCGEM2018}:
\begin{enumerate}
  \item Direct $M_\mathrm{w}$ from global moment-tensor \catalogs\ (GCMT, USGS
        W-phase, EMSC MT, regional CMT solutions).
  \item $M_\mathrm{w}$ proxies from re-computed $M_\mathrm{s}$ using an
        \textbf{alpha-trimmed median} (removing the top and bottom 10\,\% of
        contributing station estimates) to suppress outliers from poorly
        calibrated stations.
  \item $M_\mathrm{w}$ proxies from re-computed $m_\mathrm{b}$.
  \item \textbf{Exponential (non-linear) conversion models} for $M_\mathrm{s}
        \to M_\mathrm{w}$ and $m_\mathrm{b} \to M_\mathrm{w}$, preferred over
        linear relationships at magnitude extremes.
\end{enumerate}

% -----------------------------------------------------------------------------
\subsection{Magnitude Homogenisation}
\label{sec:maghomog}

\subsubsection{The target scale}

Moment magnitude $M_\mathrm{w}$ is the universal target for homogenised
\catalogs.  It is physically grounded (proportional to seismic moment $M_0$
via $M_\mathrm{w} = \tfrac{2}{3}\log_{10} M_0 - 10.7$), does not saturate at
large magnitudes, and is directly comparable across regions and recording eras
\citep{HanksKanamori1979}.

\subsubsection{Regression method: orthogonal vs.\ ordinary least squares}

Standard Ordinary Least Squares (OLS) regression is inappropriate for magnitude
conversion because \emph{both} the predictor (e.g., $M_\mathrm{s}$) and the
response ($M_\mathrm{w}$) carry measurement error.  OLS assumes the predictor
is error-free; violation of this assumption produces an attenuation bias in the
estimated slope, which in turn biases the frequency-magnitude distribution of
the converted \catalog.

The current consensus best practice is \textbf{Generalised Orthogonal
Regression} (GOR), sometimes called errors-in-variables regression or
orthogonal distance regression \citep{Grunthal2016, Akkar2014, Cotton2013,
DeadSea2022}.  GOR minimises the sum of squared perpendicular distances from
the data points to the regression line, weighted by the ratio of the variances
of the two variables:

\begin{equation}
  \eta = \frac{\sigma^2_{M_\mathrm{w}}}{\sigma^2_{M_x}},
  \label{eq:eta}
\end{equation}

where $M_x$ is the predictor magnitude type.  When $\eta = 1$ (equal
uncertainties), GOR reduces to the standard geometric mean regression.

\paragraph{Functional form.}
Linear GOR is adequate for most of the magnitude range.  The ISC-GEM pipeline
explicitly chose \emph{exponential (non-linear) relationships} for
$M_\mathrm{s} \to M_\mathrm{w}$ and $m_\mathrm{b} \to M_\mathrm{w}$, citing
better behaviour at extremes where both scales have physical non-linearities
\citep{ISCGEM2018}.  Piecewise-linear models with optimised crossover points
(as in the CPTI15 Italian \catalog\ pipeline, \citealt{CPTI15}) offer a
practical middle ground.

The commonly cited \citet{Scordilis2006} linear coefficients
(e.g., $M_\mathrm{w} = 0.67\,M_\mathrm{L} + 1.17$) are derived from OLS and
should be understood as an approximation; they systematically overestimate
$M_\mathrm{w}$ for small events and underestimate it near the saturation
boundary of each scale.

\subsubsection{Magnitude type priority hierarchy}

When multiple magnitude estimates of different types coexist for one event, the
priority order recommended by ISC-GEM and supported by the broader literature
\citep{ISCGEM2018, Grunthal2016} is:

\begin{enumerate}
  \item $M_\mathrm{w}$ from GCMT waveform inversion (gold standard)
  \item $M_\mathrm{w}$ from USGS/NEIC W-phase or broadband solution
  \item $M_\mathrm{w}$ from regional CMT (e.g., GeoNet CMT, EMSC MT)
  \item $M_\mathrm{w}$ proxy converted from $M_\mathrm{s}$
  \item $M_\mathrm{w}$ proxy converted from $m_\mathrm{b}$
  \item $M_\mathrm{w}$ proxy converted from $M_\mathrm{L}$
  \item $M_\mathrm{w}$ proxy converted from $M_\mathrm{d}$ (least reliable)
\end{enumerate}

Within each priority tier, the estimate with the lower reported uncertainty is
preferred \citep{ISCGEM2018}.

\subsubsection{Bayesian magnitude merging}

When several authoritative $M_\mathrm{w}$ estimates are available for the same
event, simple priority selection discards valid information.  \citet{Schorlemmer2024}
proposed a Bayesian framework in which each network's observed magnitude is
modelled as a linear function of the unknown true magnitude:

\begin{equation}
  M_{\mathrm{obs},i} = a_i + b_i\,M_{\mathrm{true}} + \varepsilon_i,
  \qquad \varepsilon_i \sim \mathcal{N}(0,\,\sigma_i^2),
  \label{eq:bayesmag}
\end{equation}

with a Gutenberg-Richter prior on $M_{\mathrm{true}}$.  The posterior mean
provides an optimal merged estimate.  When two authoritative solutions agree
within $0.1\,M$ units, the posterior is insensitive to the choice of prior;
when they disagree by more than $0.3\,M$ units, a data quality flag should be
raised regardless of which method is used \citep{Schorlemmer2024}.

\paragraph{Uncertainty propagation.}
For events where the conversion is applied, the propagated uncertainty is

\begin{equation}
  \sigma_{M_\mathrm{w}}^2
      = \sigma_{\hat{a}}^2
      + \left(\frac{\partial f}{\partial M_x}\right)^2 \sigma_{M_x}^2,
  \label{eq:uncert}
\end{equation}

where $\hat{a}$ represents the conversion-equation intercept uncertainty and
$f(M_x)$ is the conversion function.  This formulation, used in
\citet{DeadSea2022} and \citet{Grunthal2016}, ensures that uncertainty is
correctly amplified when the conversion slope departs from unity (e.g., for
$m_\mathrm{b}$ where $\partial f/\partial m_b \approx 0.85$).

\subsubsection{Combining macroseismic and instrumental estimates}

Where both macroseismic ($M_{\mathrm{w,macro}}$) and instrumental
($M_{\mathrm{w,inst}}$) estimates exist, the CPTI15 pipeline
\citep{CPTI15} uses inverse-variance weighting:

\begin{equation}
  M_{\mathrm{w,comb}}
    = \frac{
        M_{\mathrm{w,macro}} / \sigma_{\mathrm{macro}}^2
        + M_{\mathrm{w,inst}} / \sigma_{\mathrm{inst}}^2
      }{
        1/\sigma_{\mathrm{macro}}^2 + 1/\sigma_{\mathrm{inst}}^2
      }.
  \label{eq:ivw}
\end{equation}

% -----------------------------------------------------------------------------
\subsection{Network Authority Hierarchies}
\label{sec:authority}

\subsubsection{The ANSS ComCat model}

The ANSS Comprehensive Earthquake \catalog\ assigns a \texttt{preferredWeight}
to each product submission from each contributing network \citep{usgs_comcat}.
The network with the highest weight for a given product type (origin, magnitude,
focal mechanism) is the preferred source.  Where weights are equal, the most
recently updated product takes precedence.  Regional networks are authoritative
within their monitoring footprint; USGS NEIC fills gaps globally
\citep{usgs_comcat, Woessner2005}.

\subsubsection{The ISC model}

The ISC applies a fundamentally different philosophy: it does not simply pass
through one agency's solution.  Instead it re-locates \emph{every} event using
all contributed phase data and designates the ISC-authored hypocenter as prime.
All agency solutions remain in the bulletin as secondary hypocenters
\citep{ISCBulletinRebuild2020}.  This preserves full traceability and is the
gold standard for research-grade \catalogs.

\subsubsection{A recommended default hierarchy}

Drawing on the ANSS, ISC-GEM, and EMSC models, \cref{tab:hierarchy} lists a
practical global priority ranking.  Regional events should override the global
ranking when a specialist regional network has lower azimuthal gap and higher
station count than any global agency for the event in question.

\begin{table}[H]
  \centering
  \caption{Recommended global network authority hierarchy for parameter
           selection.  Priority 1 = most authoritative.  ``Regional override''
           indicates that this network takes precedence for events within its
           defined geographic footprint.}
  \label{tab:hierarchy}
  \begin{tabular}{cllp{4.5cm}}
    \toprule
    Priority & Network / Agency & Code & Notes \\
    \midrule
    1 & GCMT / Global CMT & \texttt{gcmt} & Waveform $M_\mathrm{w}$; gold standard \\
    2 & ISC / ISC-GEM & \texttt{isc} & Relocated; full phase data \\
    3 & USGS NEIC & \texttt{us} & Global; W-phase $M_\mathrm{w}$ \\
    4 & EMSC & \texttt{emsc} & Euro-Med aggregator \\
    5 & JMA & \texttt{jp} & Regional override: Japan \\
    6 & GeoNet / GNS & \texttt{nz} & Regional override: Aotearoa NZ \\
    7 & GEOFON / GFZ & \texttt{gfz} & Broadband; rapid solutions \\
    8 & INGV & \texttt{ingv} & Regional override: Italy \\
    9 & IGN & \texttt{ign} & Regional override: Spain \\
    \bottomrule
  \end{tabular}
\end{table}

% -----------------------------------------------------------------------------
\subsection{Parameter Selection Within a Merged Group}
\label{sec:selection}

Once all records for a physical earthquake are grouped (via any of the methods
in \cref{sec:matching}), the following selection rules reflect the consensus
of the ISC-GEM, ANSS, and Dead Sea Transform pipeline documentations.

\subsubsection{Hypocenter and origin time}
Select the hypocenter from the highest-priority agency whose solution satisfies
\emph{all} of:
\begin{itemize}
  \item azimuthal gap $< 180^\circ$,
  \item at least 6 defining phases,
  \item depth not fixed (free-floating solution preferred),
  \item RMS travel-time residual $< 2.0$\,s.
\end{itemize}
The origin time is taken from the selected hypocenter and is \emph{not}
averaged across agencies \citep{Storchak2013}.

Where no agency meets all criteria, the priority order in \cref{tab:hierarchy}
determines the fallback, with a QC flag attached.

\subsubsection{Depth}
Among qualifying hypocenters, prefer the solution with the smallest depth
uncertainty.  Where uncertainty information is absent, prefer the solution with
the greater station count.  Depth solutions fixed at standard values
(e.g., 10\,km, 33\,km, 100\,km) are given lowest priority \citep{Engdahl2020}.

\subsubsection{Magnitude}
Apply the priority hierarchy in \cref{sec:maghomog}.  When multiple estimates
exist at the same priority tier, the Bayesian posterior mean of
\cref{eq:bayesmag} is preferred; the alpha-trimmed median of \citet{ISCGEM2018}
is an acceptable substitute requiring no prior specification.

\subsubsection{Focal mechanism}
Select the focal mechanism following \citet{ISCGEM2018} source authority:
GCMT $>$ regional CMT $>$ USGS first-motion $>$ automated solutions.  Within
each tier, rank by quality score integrating station polarity count, misfit, and
variance reduction.  Retain all available mechanisms as supplementary products.

% -----------------------------------------------------------------------------
\subsection{Conflict Detection and Quality Control}
\label{sec:qc}

\subsubsection{Within-group consistency flags}

\Cref{tab:qcflags} lists the primary quality flags recommended following the
ISC bulletin and ANSS ComCat quality models.  These should be computed and
stored for every merged event, not used as hard-rejection criteria.

\begin{table}[H]
  \centering
  \caption{Recommended within-group QC flags.  All flags should be stored as
           event metadata; only the most severe (``error'' severity) justify
           dissolving a group without merging.}
  \label{tab:qcflags}
  \begin{tabular}{p{3.8cm}p{2.5cm}lp{3.8cm}}
    \toprule
    Condition & Threshold & Severity & Reference \\
    \midrule
    Azimuthal gap & $> 180^\circ$ & Warning & \citet{Bondár2011} \\
    Defining phase count & $< 6$ & Warning & \citet{Bondár2011} \\
    Depth type & Fixed & Info & \citet{Engdahl2020} \\
    RMS residual & $> 2.0$\,s & Warning & \citet{ISCBulletinRebuild2020} \\
    $|\Delta M_\mathrm{w}|$ between two authoritative sources & $> 0.3$ & Warning & \citet{Schorlemmer2024} \\
    $|\Delta M_\mathrm{w}|$ converted vs.\ observed & $> 0.5$ & Warning & \citet{Grunthal2016} \\
    Group size (events from different agencies) & $> 10$ & Error & \citet{Tanaka2022} \\
    Same agency appears twice in group & any & Error & Section~\ref{sec:sameagency} \\
    Spatial spread of epicentres & $>100$\,km for $M<5$ & Warning & \citet{Benz2019} \\
    \bottomrule
  \end{tabular}
\end{table}

\subsubsection{Same-agency duplicate flag}
\label{sec:sameagency}

If the same network contributes two records to a single merged group, the group
almost certainly spans two distinct earthquakes (e.g., a foreshock-mainshock
pair or a genuine aftershock).  This should be treated as an error-severity
conflict and the group dissolved into sub-groups via a secondary pairwise
matching pass \citep{Tanaka2022, Benz2019}.

\subsubsection{Magnitude completeness audit}
\label{sec:mc}

Following \citet{Wiemer2000}, the magnitude of completeness $M_c$ should be
estimated for the merged \catalog\ in rolling temporal bins (5-year windows are
standard) using either the Maximum Curvature (MAXC) or Goodness-of-Fit (GFT)
method.  Events below $M_c$ must be flagged as potentially incomplete in any
time bin where $M_c$ is elevated, and excluded from $b$-value or seismicity
rate calculations that assume completeness \citep{Aki1965, Woessner2005}.
Sudden changes in $M_c$ should be documented alongside the network history that
explains them.

% -----------------------------------------------------------------------------
\subsection{Algorithm Recommendations and Current Implementation}
\label{sec:recommendations}

\Cref{tab:recommendations} maps the key algorithmic decisions to the methods
recommended by the literature and summarises the current state of the
implementation described in this work.

\begin{table}[H]
  \centering
  \caption{Algorithmic decisions in earthquake \catalog\ merging, recommended
           approaches from the literature, and the current implementation
           status.  GOR = Generalised Orthogonal Regression.  OLS = Ordinary
           Least Squares.}
  \label{tab:recommendations}
  \renewcommand{\arraystretch}{1.25}
  \begin{tabular}{p{3.0cm}p{3.5cm}p{4.0cm}p{2.0cm}}
    \toprule
    Decision area & Literature recommendation & Current implementation & Priority \\
    \midrule
    Event association & Graph connected-components \citep{Benz2019} & Greedy forward-sweep & Medium \\
    Dense sequences & Nearest-neighbour proximity \citep{Tanaka2022} & Adaptive time/distance windows & Low–Medium \\
    Magnitude regression & GOR / non-linear \citep{Grunthal2016, ISCGEM2018} & Linear Scordilis OLS \citep{Scordilis2006} & High \\
    Multi-estimate averaging & Alpha-trimmed median \citep{ISCGEM2018} or Bayesian \citep{Schorlemmer2024} & Priority selection & Medium \\
    Source preservation & All solutions retained (ISC model) & JSON in \texttt{source\_events} field & Low \\
    Depth selection & Uncertainty-weighted \citep{Engdahl2020} & Uncertainty-weighted & --- (done) \\
    Focal mechanism & GCMT $>$ regional CMT priority & Implemented per hierarchy & --- (done) \\
    QC flags & 9-flag ISC/ANSS model & 7 flags implemented & Low \\
    Magnitude completeness & Rolling $M_c$ in 5-year bins \citep{Wiemer2000} & Not implemented & Medium \\
    \bottomrule
  \end{tabular}
\end{table}

The highest-priority gap is the regression methodology.  Replacing the current
linear OLS Scordilis (\citeyear{Scordilis2006}) coefficients with GOR-derived
relationships—either the published SHEEC coefficients of
\citet{Grunthal2016} or region-specific coefficients computed from the
overlap period of the source \catalogs—would reduce systematic bias in the
converted frequency-magnitude distribution without requiring any architectural
changes to the merging pipeline.

The second highest-priority gap is the adoption of a graph-based connected-
components step (see \cref{sec:graph}).  The current greedy forward-sweep misses
transitive associations and processes each event against only future events in
time order, meaning that when event B has already been consumed into a group
containing A, C cannot subsequently be merged with the A-B group even if C
matches B within threshold.  A Union-Find graph implementation would resolve
this in $O(n\,\alpha(n))$ time with no change to the downstream selection logic.

% Uses natbib \citep / \citet, booktabs, amsmath, and cleveref already loaded
% in main.tex.
% =============================================================================

\section{Strategies for Merging Earthquake \catalogs}
\label{sec:merging}

Combining records from multiple seismic networks into a single unified \catalog\
is a prerequisite for seismic-hazard analysis, time-series studies of regional
seismicity, and magnitude-frequency assessments that draw on heterogeneous data
archives \citep{woessner2010instrumental, corssa_woessner}.  The task
decomposes into two conceptually distinct steps that must be treated
separately in any robust implementation.

\begin{enumerate}
  \item \textbf{Event association and deduplication.}  Determining which records
        across source \catalogs\ describe the \emph{same physical earthquake}
        and grouping them into clusters.

  \item \textbf{Parameter selection and homogenisation.}  From each cluster,
        selecting or computing a single best origin time, hypocenter, and
        magnitude, and converting all magnitude estimates to a common scale.
\end{enumerate}

Conflating these steps—for instance by applying a magnitude conversion formula
\emph{before} deciding whether two records are duplicates—introduces systematic
biases that propagate into downstream analyses
\citep{Storchak2013, Bondár2011}.

% -----------------------------------------------------------------------------
\subsection{Event Association: Spatio-Temporal Matching}
\label{sec:matching}

\subsubsection{Window-based forward-sweep (classical approach)}

The oldest and most widely deployed method sorts all incoming origin reports by
origin time and scans forward with a sliding window.  If two reports from
\emph{different} networks fall within a fixed time window $\Delta t_{\max}$ and
within a fixed epicentral distance $\Delta d_{\max}$, they are flagged as
candidate duplicates \citep{NCEDCComCat}.

\paragraph{Operational thresholds.}
\Cref{tab:windows} summarises documented threshold choices from major
operational systems.

\begin{table}[H]
  \centering
  \caption{Spatio-temporal matching thresholds documented in major operational
           catalogue systems.  $\Delta t$ = time window, $\Delta d$ = distance
           window.  ANSS thresholds apply to inter-network pairs only;
           intra-network deduplication is assumed upstream.}
  \label{tab:windows}
  \begin{tabular}{llll}
    \toprule
    System & $\Delta t$ (s) & $\Delta d$ (km) & Reference \\
    \midrule
    ANSS Composite \catalog & 16 & 100 & \citet{NCEDCComCat} \\
    Typical regional build ($M \geq 4$) & 60 & 100 & \citet{Storchak2013} \\
    Local / small-event ($M < 4$) & 30 & 30 & \citet{Woessner2005} \\
    Pre-digital era (pre-1960) & 300 & 200 & \citet{Bondár2011} \\
    Dead Sea Transform region & 60 & 100 & \citet{DeadSea2022} \\
    Mexico unified \catalog & 60 & $\approx$111 ($1^{\circ}$) & \citet{Zuniga2019} \\
    \bottomrule
  \end{tabular}
\end{table}

The 16\,s / 100\,km rule of the ANSS Composite \catalog\ is the most precisely
documented public threshold for an active operational system
\citep{NCEDCComCat}.  For global teleseismic events ($M \geq 4.5$) the broader
60\,s / 100\,km window is standard; the ISC bulletin uses a comparable
proximity criterion internally when auto-grouping agency solutions prior to
relocation \citep{ISCBulletinRebuild2020}.

\paragraph{Adaptive thresholds.}
Fixed windows perform poorly across the full magnitude range because larger
events are subject to greater timing and location uncertainty at distant
networks.  The standard remedy, used in the ISC-GEM pipeline and adopted in the
current implementation, is to scale the effective window by the event magnitude
and depth \citep{ISCGEM2018}:

\begin{equation}
  \Delta t_{\mathrm{eff}} = \Delta t_{\mathrm{config}} \cdot f_t(M),
  \qquad
  \Delta d_{\mathrm{eff}} = \Delta d_{\mathrm{config}} \cdot f_d(M) \cdot g(h),
  \label{eq:adaptive}
\end{equation}

where $f_t$, $f_d$ are magnitude multipliers (e.g., 1.0 for $M < 4$, 1.5 for
$4 \leq M < 5.5$, 2.5 for $5.5 \leq M < 7$, 4.0 for $M \geq 7$) and $g(h)$
is a depth multiplier (1.0 for $h < 100$\,km, 1.2 for $100 \leq h < 300$\,km,
1.5 for $h \geq 300$\,km), consistent with ISC-GEM practice
\citep{ISCGEM2018, Engdahl2020}.

\paragraph{Limitation: dense aftershock sequences.}
During intense aftershock sequences the time gap between consecutive legitimate
aftershocks can fall well below 60\,s.  \citet{Tanaka2022} demonstrated that in
the 24 hours following the 2011 Tōhoku $M_\mathrm{w}$\,9.0 earthquake, more
than 599 genuine aftershock pairs had inter-event times under 60\,s—the same
interval used as a duplicate window.  Fixed-window methods therefore cannot
distinguish a duplicate pair from an aftershock pair in such sequences without
additional information.

% - - - - - - - - - - - - - - - - - - - - - - - - - - - - - - - - - - - - - -
\subsubsection{Nearest-neighbour proximity function}
\label{sec:nn}

\citet{Tanaka2022} introduced a physics-based proximity function that
discriminates duplicates from aftershocks by exploiting their different
topological signatures in event-pair space: duplicates form \emph{pairs}
(reports of one earthquake), whereas aftershocks form \emph{branching trees}
(one parent, many offspring).

The proximity between two events $i$ and $j$ is defined as

\begin{equation}
  \eta(i,j) = t_{ij}\; r_{ij}^{d_f}\; 10^{-b m_i},
  \label{eq:proximity}
\end{equation}

where $t_{ij}$ is the inter-event time, $r_{ij}$ is the epicentral distance,
$d_f \approx 1.6$ is the fractal dimension of the epicentre distribution,
$b \approx 1.0$ is the Gutenberg-Richter $b$-value, and $m_i$ is the magnitude
of the preceding event.  Events whose nearest-neighbour distance in this space
falls \emph{below} a threshold derived from network location error (rather than
from earthquake-physics considerations) are classified as duplicates.  Applied
to the Tōhoku sequence, duplicate identification reliability exceeded 97\,\%
\citep{Tanaka2022}.

The method requires estimates of $d_f$ and $b$ for the target seismicity, both
of which can be obtained from the source \catalogs\ themselves via standard
routines \citep{Aki1965, Wiemer2000}.

% - - - - - - - - - - - - - - - - - - - - - - - - - - - - - - - - - - - - - -
\subsubsection{Graph-based connected-components association}
\label{sec:graph}

The methodologically strongest approach constructs an undirected association
graph in which each origin report is a node and each candidate same-event pair
is an edge \citep{Benz2019}.  Connected components of this graph—found in
$O(n\,\alpha(n))$ time via the Union-Find algorithm, where $\alpha$ is the
inverse Ackermann function—correspond to clusters of reports for the same
physical earthquake.

\paragraph{Algorithm sketch.}
\begin{enumerate}
  \item Sort all reports by origin time; build a spatial index (R-tree or
        uniform grid).
  \item For each report $E$, query the index for all reports within
        $\Delta t_{\max}$ and $\Delta d_{\max}$ from \emph{different} source
        networks.
  \item Add a graph edge for each qualifying pair $(E, F)$.
  \item Find connected components; each component is one physical earthquake.
  \item Apply a preference-selection step within each component
        (\cref{sec:selection}).
\end{enumerate}

Compared with the greedy forward sweep, the graph approach naturally handles
\emph{transitive associations}: if network A matches network B, and B matches
network C, all three are placed in the same component even if A and C do not
independently satisfy the proximity criterion.  The greedy sweep can miss such
cases when B has already been marked as processed before A is examined.

\paragraph{Integer linear programming refinement.}
\citet{Benz2019} further applied Integer Linear Programming (ILP) to the
connected-component solution to resolve cases where a component might span two
closely spaced genuine earthquakes.  The ILP minimises the number of source
events required to explain all associations subject to consistency constraints.
Applied to Chile 2010--2017, the method detected 817\,000+ earthquakes at
greater than 95\,\% accuracy on synthetic validation data.

The USGS NEIC Hydra production system uses a graph-association framework
conceptually similar to \citet{Benz2019} \citep{usgs_comcat}.

% -----------------------------------------------------------------------------
\subsection{The ISC-GEM Reference Methodology}
\label{sec:iscgem}

The ISC-GEM Global Instrumental Earthquake \catalog\ \citep{Storchak2013,
ISCGEM2018} covers 1904--2021 for $M \geq 5.5$ globally and $M \geq 5.0$ over
continental regions.  It represents the most thoroughly documented end-to-end
merging pipeline in the seismological literature and serves as the benchmark
against which other implementations are measured.

\subsubsection{Event association}
The ISC automated grouping program compares origin times and hypocentre
locations from all contributing agencies and assembles groups using a
proximity-scoring system that weights azimuthal gap, distance range, station
count, and phase count \citep{Bondár2011}.  Groups may be split or new events
created where the evidence supports distinct earthquakes.

\subsubsection{Relocation}
The ISC applies a two-tier relocation to all grouped events:
\begin{enumerate}
  \item \textbf{EHB algorithm} \citep{Engdahl1998}: improved hypocenters with
        particular attention to focal depth, using reidentified phases.
  \item \textbf{ISC location algorithm} \citep{Bondár2011}: fixes the depth
        output from EHB and reduces systematic location bias by modelling
        correlated travel-time prediction errors across the station network.
\end{enumerate}
The ISC-authored solution is designated \emph{prime} when sufficient phase data
are available.  All agency solutions are retained as secondary hypocenters,
preserving full audit traceability \citep{ISCBulletinRebuild2020}.

\subsubsection{Magnitude assembly}
ISC-GEM reports a single $M_\mathrm{w}$ per event, derived by \citep{ISCGEM2018}:
\begin{enumerate}
  \item Direct $M_\mathrm{w}$ from global moment-tensor \catalogs\ (GCMT, USGS
        W-phase, EMSC MT, regional CMT solutions).
  \item $M_\mathrm{w}$ proxies from re-computed $M_\mathrm{s}$ using an
        \textbf{alpha-trimmed median} (removing the top and bottom 10\,\% of
        contributing station estimates) to suppress outliers from poorly
        calibrated stations.
  \item $M_\mathrm{w}$ proxies from re-computed $m_\mathrm{b}$.
  \item \textbf{Exponential (non-linear) conversion models} for $M_\mathrm{s}
        \to M_\mathrm{w}$ and $m_\mathrm{b} \to M_\mathrm{w}$, preferred over
        linear relationships at magnitude extremes.
\end{enumerate}

% -----------------------------------------------------------------------------
\subsection{Magnitude Homogenisation}
\label{sec:maghomog}

\subsubsection{The target scale}

Moment magnitude $M_\mathrm{w}$ is the universal target for homogenised
\catalogs.  It is physically grounded (proportional to seismic moment $M_0$
via $M_\mathrm{w} = \tfrac{2}{3}\log_{10} M_0 - 10.7$), does not saturate at
large magnitudes, and is directly comparable across regions and recording eras
\citep{HanksKanamori1979}.

\subsubsection{Regression method: orthogonal vs.\ ordinary least squares}

Standard Ordinary Least Squares (OLS) regression is inappropriate for magnitude
conversion because \emph{both} the predictor (e.g., $M_\mathrm{s}$) and the
response ($M_\mathrm{w}$) carry measurement error.  OLS assumes the predictor
is error-free; violation of this assumption produces an attenuation bias in the
estimated slope, which in turn biases the frequency-magnitude distribution of
the converted \catalog.

The current consensus best practice is \textbf{Generalised Orthogonal
Regression} (GOR), sometimes called errors-in-variables regression or
orthogonal distance regression \citep{Grunthal2016, Akkar2014, Cotton2013,
DeadSea2022}.  GOR minimises the sum of squared perpendicular distances from
the data points to the regression line, weighted by the ratio of the variances
of the two variables:

\begin{equation}
  \eta = \frac{\sigma^2_{M_\mathrm{w}}}{\sigma^2_{M_x}},
  \label{eq:eta}
\end{equation}

where $M_x$ is the predictor magnitude type.  When $\eta = 1$ (equal
uncertainties), GOR reduces to the standard geometric mean regression.

\paragraph{Functional form.}
Linear GOR is adequate for most of the magnitude range.  The ISC-GEM pipeline
explicitly chose \emph{exponential (non-linear) relationships} for
$M_\mathrm{s} \to M_\mathrm{w}$ and $m_\mathrm{b} \to M_\mathrm{w}$, citing
better behaviour at extremes where both scales have physical non-linearities
\citep{ISCGEM2018}.  Piecewise-linear models with optimised crossover points
(as in the CPTI15 Italian \catalog\ pipeline, \citealt{CPTI15}) offer a
practical middle ground.

The commonly cited \citet{Scordilis2006} linear coefficients
(e.g., $M_\mathrm{w} = 0.67\,M_\mathrm{L} + 1.17$) are derived from OLS and
should be understood as an approximation; they systematically overestimate
$M_\mathrm{w}$ for small events and underestimate it near the saturation
boundary of each scale.

\subsubsection{Magnitude type priority hierarchy}

When multiple magnitude estimates of different types coexist for one event, the
priority order recommended by ISC-GEM and supported by the broader literature
\citep{ISCGEM2018, Grunthal2016} is:

\begin{enumerate}
  \item $M_\mathrm{w}$ from GCMT waveform inversion (gold standard)
  \item $M_\mathrm{w}$ from USGS/NEIC W-phase or broadband solution
  \item $M_\mathrm{w}$ from regional CMT (e.g., GeoNet CMT, EMSC MT)
  \item $M_\mathrm{w}$ proxy converted from $M_\mathrm{s}$
  \item $M_\mathrm{w}$ proxy converted from $m_\mathrm{b}$
  \item $M_\mathrm{w}$ proxy converted from $M_\mathrm{L}$
  \item $M_\mathrm{w}$ proxy converted from $M_\mathrm{d}$ (least reliable)
\end{enumerate}

Within each priority tier, the estimate with the lower reported uncertainty is
preferred \citep{ISCGEM2018}.

\subsubsection{Bayesian magnitude merging}

When several authoritative $M_\mathrm{w}$ estimates are available for the same
event, simple priority selection discards valid information.  \citet{Schorlemmer2024}
proposed a Bayesian framework in which each network's observed magnitude is
modelled as a linear function of the unknown true magnitude:

\begin{equation}
  M_{\mathrm{obs},i} = a_i + b_i\,M_{\mathrm{true}} + \varepsilon_i,
  \qquad \varepsilon_i \sim \mathcal{N}(0,\,\sigma_i^2),
  \label{eq:bayesmag}
\end{equation}

with a Gutenberg-Richter prior on $M_{\mathrm{true}}$.  The posterior mean
provides an optimal merged estimate.  When two authoritative solutions agree
within $0.1\,M$ units, the posterior is insensitive to the choice of prior;
when they disagree by more than $0.3\,M$ units, a data quality flag should be
raised regardless of which method is used \citep{Schorlemmer2024}.

\paragraph{Uncertainty propagation.}
For events where the conversion is applied, the propagated uncertainty is

\begin{equation}
  \sigma_{M_\mathrm{w}}^2
      = \sigma_{\hat{a}}^2
      + \left(\frac{\partial f}{\partial M_x}\right)^2 \sigma_{M_x}^2,
  \label{eq:uncert}
\end{equation}

where $\hat{a}$ represents the conversion-equation intercept uncertainty and
$f(M_x)$ is the conversion function.  This formulation, used in
\citet{DeadSea2022} and \citet{Grunthal2016}, ensures that uncertainty is
correctly amplified when the conversion slope departs from unity (e.g., for
$m_\mathrm{b}$ where $\partial f/\partial m_b \approx 0.85$).

\subsubsection{Combining macroseismic and instrumental estimates}

Where both macroseismic ($M_{\mathrm{w,macro}}$) and instrumental
($M_{\mathrm{w,inst}}$) estimates exist, the CPTI15 pipeline
\citep{CPTI15} uses inverse-variance weighting:

\begin{equation}
  M_{\mathrm{w,comb}}
    = \frac{
        M_{\mathrm{w,macro}} / \sigma_{\mathrm{macro}}^2
        + M_{\mathrm{w,inst}} / \sigma_{\mathrm{inst}}^2
      }{
        1/\sigma_{\mathrm{macro}}^2 + 1/\sigma_{\mathrm{inst}}^2
      }.
  \label{eq:ivw}
\end{equation}

% -----------------------------------------------------------------------------
\subsection{Network Authority Hierarchies}
\label{sec:authority}

\subsubsection{The ANSS ComCat model}

The ANSS Comprehensive Earthquake \catalog\ assigns a \texttt{preferredWeight}
to each product submission from each contributing network \citep{usgs_comcat}.
The network with the highest weight for a given product type (origin, magnitude,
focal mechanism) is the preferred source.  Where weights are equal, the most
recently updated product takes precedence.  Regional networks are authoritative
within their monitoring footprint; USGS NEIC fills gaps globally
\citep{usgs_comcat, Woessner2005}.

\subsubsection{The ISC model}

The ISC applies a fundamentally different philosophy: it does not simply pass
through one agency's solution.  Instead it re-locates \emph{every} event using
all contributed phase data and designates the ISC-authored hypocenter as prime.
All agency solutions remain in the bulletin as secondary hypocenters
\citep{ISCBulletinRebuild2020}.  This preserves full traceability and is the
gold standard for research-grade \catalogs.

\subsubsection{A recommended default hierarchy}

Drawing on the ANSS, ISC-GEM, and EMSC models, \cref{tab:hierarchy} lists a
practical global priority ranking.  Regional events should override the global
ranking when a specialist regional network has lower azimuthal gap and higher
station count than any global agency for the event in question.

\begin{table}[H]
  \centering
  \caption{Recommended global network authority hierarchy for parameter
           selection.  Priority 1 = most authoritative.  ``Regional override''
           indicates that this network takes precedence for events within its
           defined geographic footprint.}
  \label{tab:hierarchy}
  \begin{tabular}{cllp{4.5cm}}
    \toprule
    Priority & Network / Agency & Code & Notes \\
    \midrule
    1 & GCMT / Global CMT & \texttt{gcmt} & Waveform $M_\mathrm{w}$; gold standard \\
    2 & ISC / ISC-GEM & \texttt{isc} & Relocated; full phase data \\
    3 & USGS NEIC & \texttt{us} & Global; W-phase $M_\mathrm{w}$ \\
    4 & EMSC & \texttt{emsc} & Euro-Med aggregator \\
    5 & JMA & \texttt{jp} & Regional override: Japan \\
    6 & GeoNet / GNS & \texttt{nz} & Regional override: Aotearoa NZ \\
    7 & GEOFON / GFZ & \texttt{gfz} & Broadband; rapid solutions \\
    8 & INGV & \texttt{ingv} & Regional override: Italy \\
    9 & IGN & \texttt{ign} & Regional override: Spain \\
    \bottomrule
  \end{tabular}
\end{table}

% -----------------------------------------------------------------------------
\subsection{Parameter Selection Within a Merged Group}
\label{sec:selection}

Once all records for a physical earthquake are grouped (via any of the methods
in \cref{sec:matching}), the following selection rules reflect the consensus
of the ISC-GEM, ANSS, and Dead Sea Transform pipeline documentations.

\subsubsection{Hypocenter and origin time}
Select the hypocenter from the highest-priority agency whose solution satisfies
\emph{all} of:
\begin{itemize}
  \item azimuthal gap $< 180^\circ$,
  \item at least 6 defining phases,
  \item depth not fixed (free-floating solution preferred),
  \item RMS travel-time residual $< 2.0$\,s.
\end{itemize}
The origin time is taken from the selected hypocenter and is \emph{not}
averaged across agencies \citep{Storchak2013}.

Where no agency meets all criteria, the priority order in \cref{tab:hierarchy}
determines the fallback, with a QC flag attached.

\subsubsection{Depth}
Among qualifying hypocenters, prefer the solution with the smallest depth
uncertainty.  Where uncertainty information is absent, prefer the solution with
the greater station count.  Depth solutions fixed at standard values
(e.g., 10\,km, 33\,km, 100\,km) are given lowest priority \citep{Engdahl2020}.

\subsubsection{Magnitude}
Apply the priority hierarchy in \cref{sec:maghomog}.  When multiple estimates
exist at the same priority tier, the Bayesian posterior mean of
\cref{eq:bayesmag} is preferred; the alpha-trimmed median of \citet{ISCGEM2018}
is an acceptable substitute requiring no prior specification.

\subsubsection{Focal mechanism}
Select the focal mechanism following \citet{ISCGEM2018} source authority:
GCMT $>$ regional CMT $>$ USGS first-motion $>$ automated solutions.  Within
each tier, rank by quality score integrating station polarity count, misfit, and
variance reduction.  Retain all available mechanisms as supplementary products.

% -----------------------------------------------------------------------------
\subsection{Conflict Detection and Quality Control}
\label{sec:qc}

\subsubsection{Within-group consistency flags}

\Cref{tab:qcflags} lists the primary quality flags recommended following the
ISC bulletin and ANSS ComCat quality models.  These should be computed and
stored for every merged event, not used as hard-rejection criteria.

\begin{table}[H]
  \centering
  \caption{Recommended within-group QC flags.  All flags should be stored as
           event metadata; only the most severe (``error'' severity) justify
           dissolving a group without merging.}
  \label{tab:qcflags}
  \begin{tabular}{p{3.8cm}p{2.5cm}lp{3.8cm}}
    \toprule
    Condition & Threshold & Severity & Reference \\
    \midrule
    Azimuthal gap & $> 180^\circ$ & Warning & \citet{Bondár2011} \\
    Defining phase count & $< 6$ & Warning & \citet{Bondár2011} \\
    Depth type & Fixed & Info & \citet{Engdahl2020} \\
    RMS residual & $> 2.0$\,s & Warning & \citet{ISCBulletinRebuild2020} \\
    $|\Delta M_\mathrm{w}|$ between two authoritative sources & $> 0.3$ & Warning & \citet{Schorlemmer2024} \\
    $|\Delta M_\mathrm{w}|$ converted vs.\ observed & $> 0.5$ & Warning & \citet{Grunthal2016} \\
    Group size (events from different agencies) & $> 10$ & Error & \citet{Tanaka2022} \\
    Same agency appears twice in group & any & Error & Section~\ref{sec:sameagency} \\
    Spatial spread of epicentres & $>100$\,km for $M<5$ & Warning & \citet{Benz2019} \\
    \bottomrule
  \end{tabular}
\end{table}

\subsubsection{Same-agency duplicate flag}
\label{sec:sameagency}

If the same network contributes two records to a single merged group, the group
almost certainly spans two distinct earthquakes (e.g., a foreshock-mainshock
pair or a genuine aftershock).  This should be treated as an error-severity
conflict and the group dissolved into sub-groups via a secondary pairwise
matching pass \citep{Tanaka2022, Benz2019}.

\subsubsection{Magnitude completeness audit}
\label{sec:mc}

Following \citet{Wiemer2000}, the magnitude of completeness $M_c$ should be
estimated for the merged \catalog\ in rolling temporal bins (5-year windows are
standard) using either the Maximum Curvature (MAXC) or Goodness-of-Fit (GFT)
method.  Events below $M_c$ must be flagged as potentially incomplete in any
time bin where $M_c$ is elevated, and excluded from $b$-value or seismicity
rate calculations that assume completeness \citep{Aki1965, Woessner2005}.
Sudden changes in $M_c$ should be documented alongside the network history that
explains them.

% -----------------------------------------------------------------------------
\subsection{Algorithm Recommendations and Current Implementation}
\label{sec:recommendations}

\Cref{tab:recommendations} maps the key algorithmic decisions to the methods
recommended by the literature and summarises the current state of the
implementation described in this work.

\begin{table}[H]
  \centering
  \caption{Algorithmic decisions in earthquake \catalog\ merging, recommended
           approaches from the literature, and the current implementation
           status.  GOR = Generalised Orthogonal Regression.  OLS = Ordinary
           Least Squares.}
  \label{tab:recommendations}
  \renewcommand{\arraystretch}{1.25}
  \begin{tabular}{p{3.0cm}p{3.5cm}p{4.0cm}p{2.0cm}}
    \toprule
    Decision area & Literature recommendation & Current implementation & Priority \\
    \midrule
    Event association & Graph connected-components \citep{Benz2019} & Greedy forward-sweep & Medium \\
    Dense sequences & Nearest-neighbour proximity \citep{Tanaka2022} & Adaptive time/distance windows & Low–Medium \\
    Magnitude regression & GOR / non-linear \citep{Grunthal2016, ISCGEM2018} & Linear Scordilis OLS \citep{Scordilis2006} & High \\
    Multi-estimate averaging & Alpha-trimmed median \citep{ISCGEM2018} or Bayesian \citep{Schorlemmer2024} & Priority selection & Medium \\
    Source preservation & All solutions retained (ISC model) & JSON in \texttt{source\_events} field & Low \\
    Depth selection & Uncertainty-weighted \citep{Engdahl2020} & Uncertainty-weighted & --- (done) \\
    Focal mechanism & GCMT $>$ regional CMT priority & Implemented per hierarchy & --- (done) \\
    QC flags & 9-flag ISC/ANSS model & 7 flags implemented & Low \\
    Magnitude completeness & Rolling $M_c$ in 5-year bins \citep{Wiemer2000} & Not implemented & Medium \\
    \bottomrule
  \end{tabular}
\end{table}

The highest-priority gap is the regression methodology.  Replacing the current
linear OLS Scordilis (\citeyear{Scordilis2006}) coefficients with GOR-derived
relationships—either the published SHEEC coefficients of
\citet{Grunthal2016} or region-specific coefficients computed from the
overlap period of the source \catalogs—would reduce systematic bias in the
converted frequency-magnitude distribution without requiring any architectural
changes to the merging pipeline.

The second highest-priority gap is the adoption of a graph-based connected-
components step (see \cref{sec:graph}).  The current greedy forward-sweep misses
transitive associations and processes each event against only future events in
time order, meaning that when event B has already been consumed into a group
containing A, C cannot subsequently be merged with the A-B group even if C
matches B within threshold.  A Union-Find graph implementation would resolve
this in $O(n\,\alpha(n))$ time with no change to the downstream selection logic.

% Uses natbib \citep / \citet, booktabs, amsmath, and cleveref already loaded
% in main.tex.
% =============================================================================

\section{Strategies for Merging Earthquake \catalogs}
\label{sec:merging}

Combining records from multiple seismic networks into a single unified \catalog\
is a prerequisite for seismic-hazard analysis, time-series studies of regional
seismicity, and magnitude-frequency assessments that draw on heterogeneous data
archives \citep{woessner2010instrumental, corssa_woessner}.  The task
decomposes into two conceptually distinct steps that must be treated
separately in any robust implementation.

\begin{enumerate}
  \item \textbf{Event association and deduplication.}  Determining which records
        across source \catalogs\ describe the \emph{same physical earthquake}
        and grouping them into clusters.

  \item \textbf{Parameter selection and homogenisation.}  From each cluster,
        selecting or computing a single best origin time, hypocenter, and
        magnitude, and converting all magnitude estimates to a common scale.
\end{enumerate}

Conflating these steps—for instance by applying a magnitude conversion formula
\emph{before} deciding whether two records are duplicates—introduces systematic
biases that propagate into downstream analyses
\citep{Storchak2013, Bondár2011}.

% -----------------------------------------------------------------------------
\subsection{Event Association: Spatio-Temporal Matching}
\label{sec:matching}

\subsubsection{Window-based forward-sweep (classical approach)}

The oldest and most widely deployed method sorts all incoming origin reports by
origin time and scans forward with a sliding window.  If two reports from
\emph{different} networks fall within a fixed time window $\Delta t_{\max}$ and
within a fixed epicentral distance $\Delta d_{\max}$, they are flagged as
candidate duplicates \citep{NCEDCComCat}.

\paragraph{Operational thresholds.}
\Cref{tab:windows} summarises documented threshold choices from major
operational systems.

\begin{table}[H]
  \centering
  \caption{Spatio-temporal matching thresholds documented in major operational
           catalogue systems.  $\Delta t$ = time window, $\Delta d$ = distance
           window.  ANSS thresholds apply to inter-network pairs only;
           intra-network deduplication is assumed upstream.}
  \label{tab:windows}
  \begin{tabular}{llll}
    \toprule
    System & $\Delta t$ (s) & $\Delta d$ (km) & Reference \\
    \midrule
    ANSS Composite \catalog & 16 & 100 & \citet{NCEDCComCat} \\
    Typical regional build ($M \geq 4$) & 60 & 100 & \citet{Storchak2013} \\
    Local / small-event ($M < 4$) & 30 & 30 & \citet{Woessner2005} \\
    Pre-digital era (pre-1960) & 300 & 200 & \citet{Bondár2011} \\
    Dead Sea Transform region & 60 & 100 & \citet{DeadSea2022} \\
    Mexico unified \catalog & 60 & $\approx$111 ($1^{\circ}$) & \citet{Zuniga2019} \\
    \bottomrule
  \end{tabular}
\end{table}

The 16\,s / 100\,km rule of the ANSS Composite \catalog\ is the most precisely
documented public threshold for an active operational system
\citep{NCEDCComCat}.  For global teleseismic events ($M \geq 4.5$) the broader
60\,s / 100\,km window is standard; the ISC bulletin uses a comparable
proximity criterion internally when auto-grouping agency solutions prior to
relocation \citep{ISCBulletinRebuild2020}.

\paragraph{Adaptive thresholds.}
Fixed windows perform poorly across the full magnitude range because larger
events are subject to greater timing and location uncertainty at distant
networks.  The standard remedy, used in the ISC-GEM pipeline and adopted in the
current implementation, is to scale the effective window by the event magnitude
and depth \citep{ISCGEM2018}:

\begin{equation}
  \Delta t_{\mathrm{eff}} = \Delta t_{\mathrm{config}} \cdot f_t(M),
  \qquad
  \Delta d_{\mathrm{eff}} = \Delta d_{\mathrm{config}} \cdot f_d(M) \cdot g(h),
  \label{eq:adaptive}
\end{equation}

where $f_t$, $f_d$ are magnitude multipliers (e.g., 1.0 for $M < 4$, 1.5 for
$4 \leq M < 5.5$, 2.5 for $5.5 \leq M < 7$, 4.0 for $M \geq 7$) and $g(h)$
is a depth multiplier (1.0 for $h < 100$\,km, 1.2 for $100 \leq h < 300$\,km,
1.5 for $h \geq 300$\,km), consistent with ISC-GEM practice
\citep{ISCGEM2018, Engdahl2020}.

\paragraph{Limitation: dense aftershock sequences.}
During intense aftershock sequences the time gap between consecutive legitimate
aftershocks can fall well below 60\,s.  \citet{Tanaka2022} demonstrated that in
the 24 hours following the 2011 Tōhoku $M_\mathrm{w}$\,9.0 earthquake, more
than 599 genuine aftershock pairs had inter-event times under 60\,s—the same
interval used as a duplicate window.  Fixed-window methods therefore cannot
distinguish a duplicate pair from an aftershock pair in such sequences without
additional information.

% - - - - - - - - - - - - - - - - - - - - - - - - - - - - - - - - - - - - - -
\subsubsection{Nearest-neighbour proximity function}
\label{sec:nn}

\citet{Tanaka2022} introduced a physics-based proximity function that
discriminates duplicates from aftershocks by exploiting their different
topological signatures in event-pair space: duplicates form \emph{pairs}
(reports of one earthquake), whereas aftershocks form \emph{branching trees}
(one parent, many offspring).

The proximity between two events $i$ and $j$ is defined as

\begin{equation}
  \eta(i,j) = t_{ij}\; r_{ij}^{d_f}\; 10^{-b m_i},
  \label{eq:proximity}
\end{equation}

where $t_{ij}$ is the inter-event time, $r_{ij}$ is the epicentral distance,
$d_f \approx 1.6$ is the fractal dimension of the epicentre distribution,
$b \approx 1.0$ is the Gutenberg-Richter $b$-value, and $m_i$ is the magnitude
of the preceding event.  Events whose nearest-neighbour distance in this space
falls \emph{below} a threshold derived from network location error (rather than
from earthquake-physics considerations) are classified as duplicates.  Applied
to the Tōhoku sequence, duplicate identification reliability exceeded 97\,\%
\citep{Tanaka2022}.

The method requires estimates of $d_f$ and $b$ for the target seismicity, both
of which can be obtained from the source \catalogs\ themselves via standard
routines \citep{Aki1965, Wiemer2000}.

% - - - - - - - - - - - - - - - - - - - - - - - - - - - - - - - - - - - - - -
\subsubsection{Graph-based connected-components association}
\label{sec:graph}

The methodologically strongest approach constructs an undirected association
graph in which each origin report is a node and each candidate same-event pair
is an edge \citep{Benz2019}.  Connected components of this graph—found in
$O(n\,\alpha(n))$ time via the Union-Find algorithm, where $\alpha$ is the
inverse Ackermann function—correspond to clusters of reports for the same
physical earthquake.

\paragraph{Algorithm sketch.}
\begin{enumerate}
  \item Sort all reports by origin time; build a spatial index (R-tree or
        uniform grid).
  \item For each report $E$, query the index for all reports within
        $\Delta t_{\max}$ and $\Delta d_{\max}$ from \emph{different} source
        networks.
  \item Add a graph edge for each qualifying pair $(E, F)$.
  \item Find connected components; each component is one physical earthquake.
  \item Apply a preference-selection step within each component
        (\cref{sec:selection}).
\end{enumerate}

Compared with the greedy forward sweep, the graph approach naturally handles
\emph{transitive associations}: if network A matches network B, and B matches
network C, all three are placed in the same component even if A and C do not
independently satisfy the proximity criterion.  The greedy sweep can miss such
cases when B has already been marked as processed before A is examined.

\paragraph{Integer linear programming refinement.}
\citet{Benz2019} further applied Integer Linear Programming (ILP) to the
connected-component solution to resolve cases where a component might span two
closely spaced genuine earthquakes.  The ILP minimises the number of source
events required to explain all associations subject to consistency constraints.
Applied to Chile 2010--2017, the method detected 817\,000+ earthquakes at
greater than 95\,\% accuracy on synthetic validation data.

The USGS NEIC Hydra production system uses a graph-association framework
conceptually similar to \citet{Benz2019} \citep{usgs_comcat}.

% -----------------------------------------------------------------------------
\subsection{The ISC-GEM Reference Methodology}
\label{sec:iscgem}

The ISC-GEM Global Instrumental Earthquake \catalog\ \citep{Storchak2013,
ISCGEM2018} covers 1904--2021 for $M \geq 5.5$ globally and $M \geq 5.0$ over
continental regions.  It represents the most thoroughly documented end-to-end
merging pipeline in the seismological literature and serves as the benchmark
against which other implementations are measured.

\subsubsection{Event association}
The ISC automated grouping program compares origin times and hypocentre
locations from all contributing agencies and assembles groups using a
proximity-scoring system that weights azimuthal gap, distance range, station
count, and phase count \citep{Bondár2011}.  Groups may be split or new events
created where the evidence supports distinct earthquakes.

\subsubsection{Relocation}
The ISC applies a two-tier relocation to all grouped events:
\begin{enumerate}
  \item \textbf{EHB algorithm} \citep{Engdahl1998}: improved hypocenters with
        particular attention to focal depth, using reidentified phases.
  \item \textbf{ISC location algorithm} \citep{Bondár2011}: fixes the depth
        output from EHB and reduces systematic location bias by modelling
        correlated travel-time prediction errors across the station network.
\end{enumerate}
The ISC-authored solution is designated \emph{prime} when sufficient phase data
are available.  All agency solutions are retained as secondary hypocenters,
preserving full audit traceability \citep{ISCBulletinRebuild2020}.

\subsubsection{Magnitude assembly}
ISC-GEM reports a single $M_\mathrm{w}$ per event, derived by \citep{ISCGEM2018}:
\begin{enumerate}
  \item Direct $M_\mathrm{w}$ from global moment-tensor \catalogs\ (GCMT, USGS
        W-phase, EMSC MT, regional CMT solutions).
  \item $M_\mathrm{w}$ proxies from re-computed $M_\mathrm{s}$ using an
        \textbf{alpha-trimmed median} (removing the top and bottom 10\,\% of
        contributing station estimates) to suppress outliers from poorly
        calibrated stations.
  \item $M_\mathrm{w}$ proxies from re-computed $m_\mathrm{b}$.
  \item \textbf{Exponential (non-linear) conversion models} for $M_\mathrm{s}
        \to M_\mathrm{w}$ and $m_\mathrm{b} \to M_\mathrm{w}$, preferred over
        linear relationships at magnitude extremes.
\end{enumerate}

% -----------------------------------------------------------------------------
\subsection{Magnitude Homogenisation}
\label{sec:maghomog}

\subsubsection{The target scale}

Moment magnitude $M_\mathrm{w}$ is the universal target for homogenised
\catalogs.  It is physically grounded (proportional to seismic moment $M_0$
via $M_\mathrm{w} = \tfrac{2}{3}\log_{10} M_0 - 10.7$), does not saturate at
large magnitudes, and is directly comparable across regions and recording eras
\citep{HanksKanamori1979}.

\subsubsection{Regression method: orthogonal vs.\ ordinary least squares}

Standard Ordinary Least Squares (OLS) regression is inappropriate for magnitude
conversion because \emph{both} the predictor (e.g., $M_\mathrm{s}$) and the
response ($M_\mathrm{w}$) carry measurement error.  OLS assumes the predictor
is error-free; violation of this assumption produces an attenuation bias in the
estimated slope, which in turn biases the frequency-magnitude distribution of
the converted \catalog.

The current consensus best practice is \textbf{Generalised Orthogonal
Regression} (GOR), sometimes called errors-in-variables regression or
orthogonal distance regression \citep{Grunthal2016, Akkar2014, Cotton2013,
DeadSea2022}.  GOR minimises the sum of squared perpendicular distances from
the data points to the regression line, weighted by the ratio of the variances
of the two variables:

\begin{equation}
  \eta = \frac{\sigma^2_{M_\mathrm{w}}}{\sigma^2_{M_x}},
  \label{eq:eta}
\end{equation}

where $M_x$ is the predictor magnitude type.  When $\eta = 1$ (equal
uncertainties), GOR reduces to the standard geometric mean regression.

\paragraph{Functional form.}
Linear GOR is adequate for most of the magnitude range.  The ISC-GEM pipeline
explicitly chose \emph{exponential (non-linear) relationships} for
$M_\mathrm{s} \to M_\mathrm{w}$ and $m_\mathrm{b} \to M_\mathrm{w}$, citing
better behaviour at extremes where both scales have physical non-linearities
\citep{ISCGEM2018}.  Piecewise-linear models with optimised crossover points
(as in the CPTI15 Italian \catalog\ pipeline, \citealt{CPTI15}) offer a
practical middle ground.

The commonly cited \citet{Scordilis2006} linear coefficients
(e.g., $M_\mathrm{w} = 0.67\,M_\mathrm{L} + 1.17$) are derived from OLS and
should be understood as an approximation; they systematically overestimate
$M_\mathrm{w}$ for small events and underestimate it near the saturation
boundary of each scale.

\subsubsection{Magnitude type priority hierarchy}

When multiple magnitude estimates of different types coexist for one event, the
priority order recommended by ISC-GEM and supported by the broader literature
\citep{ISCGEM2018, Grunthal2016} is:

\begin{enumerate}
  \item $M_\mathrm{w}$ from GCMT waveform inversion (gold standard)
  \item $M_\mathrm{w}$ from USGS/NEIC W-phase or broadband solution
  \item $M_\mathrm{w}$ from regional CMT (e.g., GeoNet CMT, EMSC MT)
  \item $M_\mathrm{w}$ proxy converted from $M_\mathrm{s}$
  \item $M_\mathrm{w}$ proxy converted from $m_\mathrm{b}$
  \item $M_\mathrm{w}$ proxy converted from $M_\mathrm{L}$
  \item $M_\mathrm{w}$ proxy converted from $M_\mathrm{d}$ (least reliable)
\end{enumerate}

Within each priority tier, the estimate with the lower reported uncertainty is
preferred \citep{ISCGEM2018}.

\subsubsection{Bayesian magnitude merging}

When several authoritative $M_\mathrm{w}$ estimates are available for the same
event, simple priority selection discards valid information.  \citet{Schorlemmer2024}
proposed a Bayesian framework in which each network's observed magnitude is
modelled as a linear function of the unknown true magnitude:

\begin{equation}
  M_{\mathrm{obs},i} = a_i + b_i\,M_{\mathrm{true}} + \varepsilon_i,
  \qquad \varepsilon_i \sim \mathcal{N}(0,\,\sigma_i^2),
  \label{eq:bayesmag}
\end{equation}

with a Gutenberg-Richter prior on $M_{\mathrm{true}}$.  The posterior mean
provides an optimal merged estimate.  When two authoritative solutions agree
within $0.1\,M$ units, the posterior is insensitive to the choice of prior;
when they disagree by more than $0.3\,M$ units, a data quality flag should be
raised regardless of which method is used \citep{Schorlemmer2024}.

\paragraph{Uncertainty propagation.}
For events where the conversion is applied, the propagated uncertainty is

\begin{equation}
  \sigma_{M_\mathrm{w}}^2
      = \sigma_{\hat{a}}^2
      + \left(\frac{\partial f}{\partial M_x}\right)^2 \sigma_{M_x}^2,
  \label{eq:uncert}
\end{equation}

where $\hat{a}$ represents the conversion-equation intercept uncertainty and
$f(M_x)$ is the conversion function.  This formulation, used in
\citet{DeadSea2022} and \citet{Grunthal2016}, ensures that uncertainty is
correctly amplified when the conversion slope departs from unity (e.g., for
$m_\mathrm{b}$ where $\partial f/\partial m_b \approx 0.85$).

\subsubsection{Combining macroseismic and instrumental estimates}

Where both macroseismic ($M_{\mathrm{w,macro}}$) and instrumental
($M_{\mathrm{w,inst}}$) estimates exist, the CPTI15 pipeline
\citep{CPTI15} uses inverse-variance weighting:

\begin{equation}
  M_{\mathrm{w,comb}}
    = \frac{
        M_{\mathrm{w,macro}} / \sigma_{\mathrm{macro}}^2
        + M_{\mathrm{w,inst}} / \sigma_{\mathrm{inst}}^2
      }{
        1/\sigma_{\mathrm{macro}}^2 + 1/\sigma_{\mathrm{inst}}^2
      }.
  \label{eq:ivw}
\end{equation}

% -----------------------------------------------------------------------------
\subsection{Network Authority Hierarchies}
\label{sec:authority}

\subsubsection{The ANSS ComCat model}

The ANSS Comprehensive Earthquake \catalog\ assigns a \texttt{preferredWeight}
to each product submission from each contributing network \citep{usgs_comcat}.
The network with the highest weight for a given product type (origin, magnitude,
focal mechanism) is the preferred source.  Where weights are equal, the most
recently updated product takes precedence.  Regional networks are authoritative
within their monitoring footprint; USGS NEIC fills gaps globally
\citep{usgs_comcat, Woessner2005}.

\subsubsection{The ISC model}

The ISC applies a fundamentally different philosophy: it does not simply pass
through one agency's solution.  Instead it re-locates \emph{every} event using
all contributed phase data and designates the ISC-authored hypocenter as prime.
All agency solutions remain in the bulletin as secondary hypocenters
\citep{ISCBulletinRebuild2020}.  This preserves full traceability and is the
gold standard for research-grade \catalogs.

\subsubsection{A recommended default hierarchy}

Drawing on the ANSS, ISC-GEM, and EMSC models, \cref{tab:hierarchy} lists a
practical global priority ranking.  Regional events should override the global
ranking when a specialist regional network has lower azimuthal gap and higher
station count than any global agency for the event in question.

\begin{table}[H]
  \centering
  \caption{Recommended global network authority hierarchy for parameter
           selection.  Priority 1 = most authoritative.  ``Regional override''
           indicates that this network takes precedence for events within its
           defined geographic footprint.}
  \label{tab:hierarchy}
  \begin{tabular}{cllp{4.5cm}}
    \toprule
    Priority & Network / Agency & Code & Notes \\
    \midrule
    1 & GCMT / Global CMT & \texttt{gcmt} & Waveform $M_\mathrm{w}$; gold standard \\
    2 & ISC / ISC-GEM & \texttt{isc} & Relocated; full phase data \\
    3 & USGS NEIC & \texttt{us} & Global; W-phase $M_\mathrm{w}$ \\
    4 & EMSC & \texttt{emsc} & Euro-Med aggregator \\
    5 & JMA & \texttt{jp} & Regional override: Japan \\
    6 & GeoNet / GNS & \texttt{nz} & Regional override: Aotearoa NZ \\
    7 & GEOFON / GFZ & \texttt{gfz} & Broadband; rapid solutions \\
    8 & INGV & \texttt{ingv} & Regional override: Italy \\
    9 & IGN & \texttt{ign} & Regional override: Spain \\
    \bottomrule
  \end{tabular}
\end{table}

% -----------------------------------------------------------------------------
\subsection{Parameter Selection Within a Merged Group}
\label{sec:selection}

Once all records for a physical earthquake are grouped (via any of the methods
in \cref{sec:matching}), the following selection rules reflect the consensus
of the ISC-GEM, ANSS, and Dead Sea Transform pipeline documentations.

\subsubsection{Hypocenter and origin time}
Select the hypocenter from the highest-priority agency whose solution satisfies
\emph{all} of:
\begin{itemize}
  \item azimuthal gap $< 180^\circ$,
  \item at least 6 defining phases,
  \item depth not fixed (free-floating solution preferred),
  \item RMS travel-time residual $< 2.0$\,s.
\end{itemize}
The origin time is taken from the selected hypocenter and is \emph{not}
averaged across agencies \citep{Storchak2013}.

Where no agency meets all criteria, the priority order in \cref{tab:hierarchy}
determines the fallback, with a QC flag attached.

\subsubsection{Depth}
Among qualifying hypocenters, prefer the solution with the smallest depth
uncertainty.  Where uncertainty information is absent, prefer the solution with
the greater station count.  Depth solutions fixed at standard values
(e.g., 10\,km, 33\,km, 100\,km) are given lowest priority \citep{Engdahl2020}.

\subsubsection{Magnitude}
Apply the priority hierarchy in \cref{sec:maghomog}.  When multiple estimates
exist at the same priority tier, the Bayesian posterior mean of
\cref{eq:bayesmag} is preferred; the alpha-trimmed median of \citet{ISCGEM2018}
is an acceptable substitute requiring no prior specification.

\subsubsection{Focal mechanism}
Select the focal mechanism following \citet{ISCGEM2018} source authority:
GCMT $>$ regional CMT $>$ USGS first-motion $>$ automated solutions.  Within
each tier, rank by quality score integrating station polarity count, misfit, and
variance reduction.  Retain all available mechanisms as supplementary products.

% -----------------------------------------------------------------------------
\subsection{Conflict Detection and Quality Control}
\label{sec:qc}

\subsubsection{Within-group consistency flags}

\Cref{tab:qcflags} lists the primary quality flags recommended following the
ISC bulletin and ANSS ComCat quality models.  These should be computed and
stored for every merged event, not used as hard-rejection criteria.

\begin{table}[H]
  \centering
  \caption{Recommended within-group QC flags.  All flags should be stored as
           event metadata; only the most severe (``error'' severity) justify
           dissolving a group without merging.}
  \label{tab:qcflags}
  \begin{tabular}{p{3.8cm}p{2.5cm}lp{3.8cm}}
    \toprule
    Condition & Threshold & Severity & Reference \\
    \midrule
    Azimuthal gap & $> 180^\circ$ & Warning & \citet{Bondár2011} \\
    Defining phase count & $< 6$ & Warning & \citet{Bondár2011} \\
    Depth type & Fixed & Info & \citet{Engdahl2020} \\
    RMS residual & $> 2.0$\,s & Warning & \citet{ISCBulletinRebuild2020} \\
    $|\Delta M_\mathrm{w}|$ between two authoritative sources & $> 0.3$ & Warning & \citet{Schorlemmer2024} \\
    $|\Delta M_\mathrm{w}|$ converted vs.\ observed & $> 0.5$ & Warning & \citet{Grunthal2016} \\
    Group size (events from different agencies) & $> 10$ & Error & \citet{Tanaka2022} \\
    Same agency appears twice in group & any & Error & Section~\ref{sec:sameagency} \\
    Spatial spread of epicentres & $>100$\,km for $M<5$ & Warning & \citet{Benz2019} \\
    \bottomrule
  \end{tabular}
\end{table}

\subsubsection{Same-agency duplicate flag}
\label{sec:sameagency}

If the same network contributes two records to a single merged group, the group
almost certainly spans two distinct earthquakes (e.g., a foreshock-mainshock
pair or a genuine aftershock).  This should be treated as an error-severity
conflict and the group dissolved into sub-groups via a secondary pairwise
matching pass \citep{Tanaka2022, Benz2019}.

\subsubsection{Magnitude completeness audit}
\label{sec:mc}

Following \citet{Wiemer2000}, the magnitude of completeness $M_c$ should be
estimated for the merged \catalog\ in rolling temporal bins (5-year windows are
standard) using either the Maximum Curvature (MAXC) or Goodness-of-Fit (GFT)
method.  Events below $M_c$ must be flagged as potentially incomplete in any
time bin where $M_c$ is elevated, and excluded from $b$-value or seismicity
rate calculations that assume completeness \citep{Aki1965, Woessner2005}.
Sudden changes in $M_c$ should be documented alongside the network history that
explains them.

% -----------------------------------------------------------------------------
\subsection{Algorithm Recommendations and Current Implementation}
\label{sec:recommendations}

\Cref{tab:recommendations} maps the key algorithmic decisions to the methods
recommended by the literature and summarises the current state of the
implementation described in this work.

\begin{table}[H]
  \centering
  \caption{Algorithmic decisions in earthquake \catalog\ merging, recommended
           approaches from the literature, and the current implementation
           status.  GOR = Generalised Orthogonal Regression.  OLS = Ordinary
           Least Squares.}
  \label{tab:recommendations}
  \renewcommand{\arraystretch}{1.25}
  \begin{tabular}{p{3.0cm}p{3.5cm}p{4.0cm}p{2.0cm}}
    \toprule
    Decision area & Literature recommendation & Current implementation & Priority \\
    \midrule
    Event association & Graph connected-components \citep{Benz2019} & Greedy forward-sweep & Medium \\
    Dense sequences & Nearest-neighbour proximity \citep{Tanaka2022} & Adaptive time/distance windows & Low–Medium \\
    Magnitude regression & GOR / non-linear \citep{Grunthal2016, ISCGEM2018} & Linear Scordilis OLS \citep{Scordilis2006} & High \\
    Multi-estimate averaging & Alpha-trimmed median \citep{ISCGEM2018} or Bayesian \citep{Schorlemmer2024} & Priority selection & Medium \\
    Source preservation & All solutions retained (ISC model) & JSON in \texttt{source\_events} field & Low \\
    Depth selection & Uncertainty-weighted \citep{Engdahl2020} & Uncertainty-weighted & --- (done) \\
    Focal mechanism & GCMT $>$ regional CMT priority & Implemented per hierarchy & --- (done) \\
    QC flags & 9-flag ISC/ANSS model & 7 flags implemented & Low \\
    Magnitude completeness & Rolling $M_c$ in 5-year bins \citep{Wiemer2000} & Not implemented & Medium \\
    \bottomrule
  \end{tabular}
\end{table}

The highest-priority gap is the regression methodology.  Replacing the current
linear OLS Scordilis (\citeyear{Scordilis2006}) coefficients with GOR-derived
relationships—either the published SHEEC coefficients of
\citet{Grunthal2016} or region-specific coefficients computed from the
overlap period of the source \catalogs—would reduce systematic bias in the
converted frequency-magnitude distribution without requiring any architectural
changes to the merging pipeline.

The second highest-priority gap is the adoption of a graph-based connected-
components step (see \cref{sec:graph}).  The current greedy forward-sweep misses
transitive associations and processes each event against only future events in
time order, meaning that when event B has already been consumed into a group
containing A, C cannot subsequently be merged with the A-B group even if C
matches B within threshold.  A Union-Find graph implementation would resolve
this in $O(n\,\alpha(n))$ time with no change to the downstream selection logic.

% Uses natbib \citep / \citet, booktabs, amsmath, and cleveref already loaded
% in main.tex.
% =============================================================================

\section{Strategies for Merging Earthquake \catalogs}
\label{sec:merging}

Combining records from multiple seismic networks into a single unified \catalog\
is a prerequisite for seismic-hazard analysis, time-series studies of regional
seismicity, and magnitude-frequency assessments that draw on heterogeneous data
archives \citep{woessner2010instrumental, corssa_woessner}.  The task
decomposes into two conceptually distinct steps that must be treated
separately in any robust implementation.

\begin{enumerate}
  \item \textbf{Event association and deduplication.}  Determining which records
        across source \catalogs\ describe the \emph{same physical earthquake}
        and grouping them into clusters.

  \item \textbf{Parameter selection and homogenisation.}  From each cluster,
        selecting or computing a single best origin time, hypocenter, and
        magnitude, and converting all magnitude estimates to a common scale.
\end{enumerate}

Conflating these steps—for instance by applying a magnitude conversion formula
\emph{before} deciding whether two records are duplicates—introduces systematic
biases that propagate into downstream analyses
\citep{Storchak2013, Bondár2011}.

% -----------------------------------------------------------------------------
\subsection{Event Association: Spatio-Temporal Matching}
\label{sec:matching}

\subsubsection{Window-based forward-sweep (classical approach)}

The oldest and most widely deployed method sorts all incoming origin reports by
origin time and scans forward with a sliding window.  If two reports from
\emph{different} networks fall within a fixed time window $\Delta t_{\max}$ and
within a fixed epicentral distance $\Delta d_{\max}$, they are flagged as
candidate duplicates \citep{NCEDCComCat}.

\paragraph{Operational thresholds.}
\Cref{tab:windows} summarises documented threshold choices from major
operational systems.

\begin{table}[H]
  \centering
  \caption{Spatio-temporal matching thresholds documented in major operational
           catalogue systems.  $\Delta t$ = time window, $\Delta d$ = distance
           window.  ANSS thresholds apply to inter-network pairs only;
           intra-network deduplication is assumed upstream.}
  \label{tab:windows}
  \begin{tabular}{llll}
    \toprule
    System & $\Delta t$ (s) & $\Delta d$ (km) & Reference \\
    \midrule
    ANSS Composite \catalog & 16 & 100 & \citet{NCEDCComCat} \\
    Typical regional build ($M \geq 4$) & 60 & 100 & \citet{Storchak2013} \\
    Local / small-event ($M < 4$) & 30 & 30 & \citet{Woessner2005} \\
    Pre-digital era (pre-1960) & 300 & 200 & \citet{Bondár2011} \\
    Dead Sea Transform region & 60 & 100 & \citet{DeadSea2022} \\
    Mexico unified \catalog & 60 & $\approx$111 ($1^{\circ}$) & \citet{Zuniga2019} \\
    \bottomrule
  \end{tabular}
\end{table}

The 16\,s / 100\,km rule of the ANSS Composite \catalog\ is the most precisely
documented public threshold for an active operational system
\citep{NCEDCComCat}.  For global teleseismic events ($M \geq 4.5$) the broader
60\,s / 100\,km window is standard; the ISC bulletin uses a comparable
proximity criterion internally when auto-grouping agency solutions prior to
relocation \citep{ISCBulletinRebuild2020}.

\paragraph{Adaptive thresholds.}
Fixed windows perform poorly across the full magnitude range because larger
events are subject to greater timing and location uncertainty at distant
networks.  The standard remedy, used in the ISC-GEM pipeline and adopted in the
current implementation, is to scale the effective window by the event magnitude
and depth \citep{ISCGEM2018}:

\begin{equation}
  \Delta t_{\mathrm{eff}} = \Delta t_{\mathrm{config}} \cdot f_t(M),
  \qquad
  \Delta d_{\mathrm{eff}} = \Delta d_{\mathrm{config}} \cdot f_d(M) \cdot g(h),
  \label{eq:adaptive}
\end{equation}

where $f_t$, $f_d$ are magnitude multipliers (e.g., 1.0 for $M < 4$, 1.5 for
$4 \leq M < 5.5$, 2.5 for $5.5 \leq M < 7$, 4.0 for $M \geq 7$) and $g(h)$
is a depth multiplier (1.0 for $h < 100$\,km, 1.2 for $100 \leq h < 300$\,km,
1.5 for $h \geq 300$\,km), consistent with ISC-GEM practice
\citep{ISCGEM2018, Engdahl2020}.

\paragraph{Limitation: dense aftershock sequences.}
During intense aftershock sequences the time gap between consecutive legitimate
aftershocks can fall well below 60\,s.  \citet{Tanaka2022} demonstrated that in
the 24 hours following the 2011 Tōhoku $M_\mathrm{w}$\,9.0 earthquake, more
than 599 genuine aftershock pairs had inter-event times under 60\,s—the same
interval used as a duplicate window.  Fixed-window methods therefore cannot
distinguish a duplicate pair from an aftershock pair in such sequences without
additional information.

% - - - - - - - - - - - - - - - - - - - - - - - - - - - - - - - - - - - - - -
\subsubsection{Nearest-neighbour proximity function}
\label{sec:nn}

\citet{Tanaka2022} introduced a physics-based proximity function that
discriminates duplicates from aftershocks by exploiting their different
topological signatures in event-pair space: duplicates form \emph{pairs}
(reports of one earthquake), whereas aftershocks form \emph{branching trees}
(one parent, many offspring).

The proximity between two events $i$ and $j$ is defined as

\begin{equation}
  \eta(i,j) = t_{ij}\; r_{ij}^{d_f}\; 10^{-b m_i},
  \label{eq:proximity}
\end{equation}

where $t_{ij}$ is the inter-event time, $r_{ij}$ is the epicentral distance,
$d_f \approx 1.6$ is the fractal dimension of the epicentre distribution,
$b \approx 1.0$ is the Gutenberg-Richter $b$-value, and $m_i$ is the magnitude
of the preceding event.  Events whose nearest-neighbour distance in this space
falls \emph{below} a threshold derived from network location error (rather than
from earthquake-physics considerations) are classified as duplicates.  Applied
to the Tōhoku sequence, duplicate identification reliability exceeded 97\,\%
\citep{Tanaka2022}.

The method requires estimates of $d_f$ and $b$ for the target seismicity, both
of which can be obtained from the source \catalogs\ themselves via standard
routines \citep{Aki1965, Wiemer2000}.

% - - - - - - - - - - - - - - - - - - - - - - - - - - - - - - - - - - - - - -
\subsubsection{Graph-based connected-components association}
\label{sec:graph}

The methodologically strongest approach constructs an undirected association
graph in which each origin report is a node and each candidate same-event pair
is an edge \citep{Benz2019}.  Connected components of this graph—found in
$O(n\,\alpha(n))$ time via the Union-Find algorithm, where $\alpha$ is the
inverse Ackermann function—correspond to clusters of reports for the same
physical earthquake.

\paragraph{Algorithm sketch.}
\begin{enumerate}
  \item Sort all reports by origin time; build a spatial index (R-tree or
        uniform grid).
  \item For each report $E$, query the index for all reports within
        $\Delta t_{\max}$ and $\Delta d_{\max}$ from \emph{different} source
        networks.
  \item Add a graph edge for each qualifying pair $(E, F)$.
  \item Find connected components; each component is one physical earthquake.
  \item Apply a preference-selection step within each component
        (\cref{sec:selection}).
\end{enumerate}

Compared with the greedy forward sweep, the graph approach naturally handles
\emph{transitive associations}: if network A matches network B, and B matches
network C, all three are placed in the same component even if A and C do not
independently satisfy the proximity criterion.  The greedy sweep can miss such
cases when B has already been marked as processed before A is examined.

\paragraph{Integer linear programming refinement.}
\citet{Benz2019} further applied Integer Linear Programming (ILP) to the
connected-component solution to resolve cases where a component might span two
closely spaced genuine earthquakes.  The ILP minimises the number of source
events required to explain all associations subject to consistency constraints.
Applied to Chile 2010--2017, the method detected 817\,000+ earthquakes at
greater than 95\,\% accuracy on synthetic validation data.

The USGS NEIC Hydra production system uses a graph-association framework
conceptually similar to \citet{Benz2019} \citep{usgs_comcat}.

% -----------------------------------------------------------------------------
\subsection{The ISC-GEM Reference Methodology}
\label{sec:iscgem}

The ISC-GEM Global Instrumental Earthquake \catalog\ \citep{Storchak2013,
ISCGEM2018} covers 1904--2021 for $M \geq 5.5$ globally and $M \geq 5.0$ over
continental regions.  It represents the most thoroughly documented end-to-end
merging pipeline in the seismological literature and serves as the benchmark
against which other implementations are measured.

\subsubsection{Event association}
The ISC automated grouping program compares origin times and hypocentre
locations from all contributing agencies and assembles groups using a
proximity-scoring system that weights azimuthal gap, distance range, station
count, and phase count \citep{Bondár2011}.  Groups may be split or new events
created where the evidence supports distinct earthquakes.

\subsubsection{Relocation}
The ISC applies a two-tier relocation to all grouped events:
\begin{enumerate}
  \item \textbf{EHB algorithm} \citep{Engdahl1998}: improved hypocenters with
        particular attention to focal depth, using reidentified phases.
  \item \textbf{ISC location algorithm} \citep{Bondár2011}: fixes the depth
        output from EHB and reduces systematic location bias by modelling
        correlated travel-time prediction errors across the station network.
\end{enumerate}
The ISC-authored solution is designated \emph{prime} when sufficient phase data
are available.  All agency solutions are retained as secondary hypocenters,
preserving full audit traceability \citep{ISCBulletinRebuild2020}.

\subsubsection{Magnitude assembly}
ISC-GEM reports a single $M_\mathrm{w}$ per event, derived by \citep{ISCGEM2018}:
\begin{enumerate}
  \item Direct $M_\mathrm{w}$ from global moment-tensor \catalogs\ (GCMT, USGS
        W-phase, EMSC MT, regional CMT solutions).
  \item $M_\mathrm{w}$ proxies from re-computed $M_\mathrm{s}$ using an
        \textbf{alpha-trimmed median} (removing the top and bottom 10\,\% of
        contributing station estimates) to suppress outliers from poorly
        calibrated stations.
  \item $M_\mathrm{w}$ proxies from re-computed $m_\mathrm{b}$.
  \item \textbf{Exponential (non-linear) conversion models} for $M_\mathrm{s}
        \to M_\mathrm{w}$ and $m_\mathrm{b} \to M_\mathrm{w}$, preferred over
        linear relationships at magnitude extremes.
\end{enumerate}

% -----------------------------------------------------------------------------
\subsection{Magnitude Homogenisation}
\label{sec:maghomog}

\subsubsection{The target scale}

Moment magnitude $M_\mathrm{w}$ is the universal target for homogenised
\catalogs.  It is physically grounded (proportional to seismic moment $M_0$
via $M_\mathrm{w} = \tfrac{2}{3}\log_{10} M_0 - 10.7$), does not saturate at
large magnitudes, and is directly comparable across regions and recording eras
\citep{HanksKanamori1979}.

\subsubsection{Regression method: orthogonal vs.\ ordinary least squares}

Standard Ordinary Least Squares (OLS) regression is inappropriate for magnitude
conversion because \emph{both} the predictor (e.g., $M_\mathrm{s}$) and the
response ($M_\mathrm{w}$) carry measurement error.  OLS assumes the predictor
is error-free; violation of this assumption produces an attenuation bias in the
estimated slope, which in turn biases the frequency-magnitude distribution of
the converted \catalog.

The current consensus best practice is \textbf{Generalised Orthogonal
Regression} (GOR), sometimes called errors-in-variables regression or
orthogonal distance regression \citep{Grunthal2016, Akkar2014, Cotton2013,
DeadSea2022}.  GOR minimises the sum of squared perpendicular distances from
the data points to the regression line, weighted by the ratio of the variances
of the two variables:

\begin{equation}
  \eta = \frac{\sigma^2_{M_\mathrm{w}}}{\sigma^2_{M_x}},
  \label{eq:eta}
\end{equation}

where $M_x$ is the predictor magnitude type.  When $\eta = 1$ (equal
uncertainties), GOR reduces to the standard geometric mean regression.

\paragraph{Functional form.}
Linear GOR is adequate for most of the magnitude range.  The ISC-GEM pipeline
explicitly chose \emph{exponential (non-linear) relationships} for
$M_\mathrm{s} \to M_\mathrm{w}$ and $m_\mathrm{b} \to M_\mathrm{w}$, citing
better behaviour at extremes where both scales have physical non-linearities
\citep{ISCGEM2018}.  Piecewise-linear models with optimised crossover points
(as in the CPTI15 Italian \catalog\ pipeline, \citealt{CPTI15}) offer a
practical middle ground.

The commonly cited \citet{Scordilis2006} linear coefficients
(e.g., $M_\mathrm{w} = 0.67\,M_\mathrm{L} + 1.17$) are derived from OLS and
should be understood as an approximation; they systematically overestimate
$M_\mathrm{w}$ for small events and underestimate it near the saturation
boundary of each scale.

\subsubsection{Magnitude type priority hierarchy}

When multiple magnitude estimates of different types coexist for one event, the
priority order recommended by ISC-GEM and supported by the broader literature
\citep{ISCGEM2018, Grunthal2016} is:

\begin{enumerate}
  \item $M_\mathrm{w}$ from GCMT waveform inversion (gold standard)
  \item $M_\mathrm{w}$ from USGS/NEIC W-phase or broadband solution
  \item $M_\mathrm{w}$ from regional CMT (e.g., GeoNet CMT, EMSC MT)
  \item $M_\mathrm{w}$ proxy converted from $M_\mathrm{s}$
  \item $M_\mathrm{w}$ proxy converted from $m_\mathrm{b}$
  \item $M_\mathrm{w}$ proxy converted from $M_\mathrm{L}$
  \item $M_\mathrm{w}$ proxy converted from $M_\mathrm{d}$ (least reliable)
\end{enumerate}

Within each priority tier, the estimate with the lower reported uncertainty is
preferred \citep{ISCGEM2018}.

\subsubsection{Bayesian magnitude merging}

When several authoritative $M_\mathrm{w}$ estimates are available for the same
event, simple priority selection discards valid information.  \citet{Schorlemmer2024}
proposed a Bayesian framework in which each network's observed magnitude is
modelled as a linear function of the unknown true magnitude:

\begin{equation}
  M_{\mathrm{obs},i} = a_i + b_i\,M_{\mathrm{true}} + \varepsilon_i,
  \qquad \varepsilon_i \sim \mathcal{N}(0,\,\sigma_i^2),
  \label{eq:bayesmag}
\end{equation}

with a Gutenberg-Richter prior on $M_{\mathrm{true}}$.  The posterior mean
provides an optimal merged estimate.  When two authoritative solutions agree
within $0.1\,M$ units, the posterior is insensitive to the choice of prior;
when they disagree by more than $0.3\,M$ units, a data quality flag should be
raised regardless of which method is used \citep{Schorlemmer2024}.

\paragraph{Uncertainty propagation.}
For events where the conversion is applied, the propagated uncertainty is

\begin{equation}
  \sigma_{M_\mathrm{w}}^2
      = \sigma_{\hat{a}}^2
      + \left(\frac{\partial f}{\partial M_x}\right)^2 \sigma_{M_x}^2,
  \label{eq:uncert}
\end{equation}

where $\hat{a}$ represents the conversion-equation intercept uncertainty and
$f(M_x)$ is the conversion function.  This formulation, used in
\citet{DeadSea2022} and \citet{Grunthal2016}, ensures that uncertainty is
correctly amplified when the conversion slope departs from unity (e.g., for
$m_\mathrm{b}$ where $\partial f/\partial m_b \approx 0.85$).

\subsubsection{Combining macroseismic and instrumental estimates}

Where both macroseismic ($M_{\mathrm{w,macro}}$) and instrumental
($M_{\mathrm{w,inst}}$) estimates exist, the CPTI15 pipeline
\citep{CPTI15} uses inverse-variance weighting:

\begin{equation}
  M_{\mathrm{w,comb}}
    = \frac{
        M_{\mathrm{w,macro}} / \sigma_{\mathrm{macro}}^2
        + M_{\mathrm{w,inst}} / \sigma_{\mathrm{inst}}^2
      }{
        1/\sigma_{\mathrm{macro}}^2 + 1/\sigma_{\mathrm{inst}}^2
      }.
  \label{eq:ivw}
\end{equation}

% -----------------------------------------------------------------------------
\subsection{Network Authority Hierarchies}
\label{sec:authority}

\subsubsection{The ANSS ComCat model}

The ANSS Comprehensive Earthquake \catalog\ assigns a \texttt{preferredWeight}
to each product submission from each contributing network \citep{usgs_comcat}.
The network with the highest weight for a given product type (origin, magnitude,
focal mechanism) is the preferred source.  Where weights are equal, the most
recently updated product takes precedence.  Regional networks are authoritative
within their monitoring footprint; USGS NEIC fills gaps globally
\citep{usgs_comcat, Woessner2005}.

\subsubsection{The ISC model}

The ISC applies a fundamentally different philosophy: it does not simply pass
through one agency's solution.  Instead it re-locates \emph{every} event using
all contributed phase data and designates the ISC-authored hypocenter as prime.
All agency solutions remain in the bulletin as secondary hypocenters
\citep{ISCBulletinRebuild2020}.  This preserves full traceability and is the
gold standard for research-grade \catalogs.

\subsubsection{A recommended default hierarchy}

Drawing on the ANSS, ISC-GEM, and EMSC models, \cref{tab:hierarchy} lists a
practical global priority ranking.  Regional events should override the global
ranking when a specialist regional network has lower azimuthal gap and higher
station count than any global agency for the event in question.

\begin{table}[H]
  \centering
  \caption{Recommended global network authority hierarchy for parameter
           selection.  Priority 1 = most authoritative.  ``Regional override''
           indicates that this network takes precedence for events within its
           defined geographic footprint.}
  \label{tab:hierarchy}
  \begin{tabular}{cllp{4.5cm}}
    \toprule
    Priority & Network / Agency & Code & Notes \\
    \midrule
    1 & GCMT / Global CMT & \texttt{gcmt} & Waveform $M_\mathrm{w}$; gold standard \\
    2 & ISC / ISC-GEM & \texttt{isc} & Relocated; full phase data \\
    3 & USGS NEIC & \texttt{us} & Global; W-phase $M_\mathrm{w}$ \\
    4 & EMSC & \texttt{emsc} & Euro-Med aggregator \\
    5 & JMA & \texttt{jp} & Regional override: Japan \\
    6 & GeoNet / GNS & \texttt{nz} & Regional override: Aotearoa NZ \\
    7 & GEOFON / GFZ & \texttt{gfz} & Broadband; rapid solutions \\
    8 & INGV & \texttt{ingv} & Regional override: Italy \\
    9 & IGN & \texttt{ign} & Regional override: Spain \\
    \bottomrule
  \end{tabular}
\end{table}

% -----------------------------------------------------------------------------
\subsection{Parameter Selection Within a Merged Group}
\label{sec:selection}

Once all records for a physical earthquake are grouped (via any of the methods
in \cref{sec:matching}), the following selection rules reflect the consensus
of the ISC-GEM, ANSS, and Dead Sea Transform pipeline documentations.

\subsubsection{Hypocenter and origin time}
Select the hypocenter from the highest-priority agency whose solution satisfies
\emph{all} of:
\begin{itemize}
  \item azimuthal gap $< 180^\circ$,
  \item at least 6 defining phases,
  \item depth not fixed (free-floating solution preferred),
  \item RMS travel-time residual $< 2.0$\,s.
\end{itemize}
The origin time is taken from the selected hypocenter and is \emph{not}
averaged across agencies \citep{Storchak2013}.

Where no agency meets all criteria, the priority order in \cref{tab:hierarchy}
determines the fallback, with a QC flag attached.

\subsubsection{Depth}
Among qualifying hypocenters, prefer the solution with the smallest depth
uncertainty.  Where uncertainty information is absent, prefer the solution with
the greater station count.  Depth solutions fixed at standard values
(e.g., 10\,km, 33\,km, 100\,km) are given lowest priority \citep{Engdahl2020}.

\subsubsection{Magnitude}
Apply the priority hierarchy in \cref{sec:maghomog}.  When multiple estimates
exist at the same priority tier, the Bayesian posterior mean of
\cref{eq:bayesmag} is preferred; the alpha-trimmed median of \citet{ISCGEM2018}
is an acceptable substitute requiring no prior specification.

\subsubsection{Focal mechanism}
Select the focal mechanism following \citet{ISCGEM2018} source authority:
GCMT $>$ regional CMT $>$ USGS first-motion $>$ automated solutions.  Within
each tier, rank by quality score integrating station polarity count, misfit, and
variance reduction.  Retain all available mechanisms as supplementary products.

% -----------------------------------------------------------------------------
\subsection{Conflict Detection and Quality Control}
\label{sec:qc}

\subsubsection{Within-group consistency flags}

\Cref{tab:qcflags} lists the primary quality flags recommended following the
ISC bulletin and ANSS ComCat quality models.  These should be computed and
stored for every merged event, not used as hard-rejection criteria.

\begin{table}[H]
  \centering
  \caption{Recommended within-group QC flags.  All flags should be stored as
           event metadata; only the most severe (``error'' severity) justify
           dissolving a group without merging.}
  \label{tab:qcflags}
  \begin{tabular}{p{3.8cm}p{2.5cm}lp{3.8cm}}
    \toprule
    Condition & Threshold & Severity & Reference \\
    \midrule
    Azimuthal gap & $> 180^\circ$ & Warning & \citet{Bondár2011} \\
    Defining phase count & $< 6$ & Warning & \citet{Bondár2011} \\
    Depth type & Fixed & Info & \citet{Engdahl2020} \\
    RMS residual & $> 2.0$\,s & Warning & \citet{ISCBulletinRebuild2020} \\
    $|\Delta M_\mathrm{w}|$ between two authoritative sources & $> 0.3$ & Warning & \citet{Schorlemmer2024} \\
    $|\Delta M_\mathrm{w}|$ converted vs.\ observed & $> 0.5$ & Warning & \citet{Grunthal2016} \\
    Group size (events from different agencies) & $> 10$ & Error & \citet{Tanaka2022} \\
    Same agency appears twice in group & any & Error & Section~\ref{sec:sameagency} \\
    Spatial spread of epicentres & $>100$\,km for $M<5$ & Warning & \citet{Benz2019} \\
    \bottomrule
  \end{tabular}
\end{table}

\subsubsection{Same-agency duplicate flag}
\label{sec:sameagency}

If the same network contributes two records to a single merged group, the group
almost certainly spans two distinct earthquakes (e.g., a foreshock-mainshock
pair or a genuine aftershock).  This should be treated as an error-severity
conflict and the group dissolved into sub-groups via a secondary pairwise
matching pass \citep{Tanaka2022, Benz2019}.

\subsubsection{Magnitude completeness audit}
\label{sec:mc}

Following \citet{Wiemer2000}, the magnitude of completeness $M_c$ should be
estimated for the merged \catalog\ in rolling temporal bins (5-year windows are
standard) using either the Maximum Curvature (MAXC) or Goodness-of-Fit (GFT)
method.  Events below $M_c$ must be flagged as potentially incomplete in any
time bin where $M_c$ is elevated, and excluded from $b$-value or seismicity
rate calculations that assume completeness \citep{Aki1965, Woessner2005}.
Sudden changes in $M_c$ should be documented alongside the network history that
explains them.

% -----------------------------------------------------------------------------
\subsection{Algorithm Recommendations and Current Implementation}
\label{sec:recommendations}

\Cref{tab:recommendations} maps the key algorithmic decisions to the methods
recommended by the literature and summarises the current state of the
implementation described in this work.

\begin{table}[H]
  \centering
  \caption{Algorithmic decisions in earthquake \catalog\ merging, recommended
           approaches from the literature, and the current implementation
           status.  GOR = Generalised Orthogonal Regression.  OLS = Ordinary
           Least Squares.}
  \label{tab:recommendations}
  \renewcommand{\arraystretch}{1.25}
  \begin{tabular}{p{3.0cm}p{3.5cm}p{4.0cm}p{2.0cm}}
    \toprule
    Decision area & Literature recommendation & Current implementation & Priority \\
    \midrule
    Event association & Graph connected-components \citep{Benz2019} & Greedy forward-sweep & Medium \\
    Dense sequences & Nearest-neighbour proximity \citep{Tanaka2022} & Adaptive time/distance windows & Low–Medium \\
    Magnitude regression & GOR / non-linear \citep{Grunthal2016, ISCGEM2018} & Linear Scordilis OLS \citep{Scordilis2006} & High \\
    Multi-estimate averaging & Alpha-trimmed median \citep{ISCGEM2018} or Bayesian \citep{Schorlemmer2024} & Priority selection & Medium \\
    Source preservation & All solutions retained (ISC model) & JSON in \texttt{source\_events} field & Low \\
    Depth selection & Uncertainty-weighted \citep{Engdahl2020} & Uncertainty-weighted & --- (done) \\
    Focal mechanism & GCMT $>$ regional CMT priority & Implemented per hierarchy & --- (done) \\
    QC flags & 9-flag ISC/ANSS model & 7 flags implemented & Low \\
    Magnitude completeness & Rolling $M_c$ in 5-year bins \citep{Wiemer2000} & Not implemented & Medium \\
    \bottomrule
  \end{tabular}
\end{table}

The highest-priority gap is the regression methodology.  Replacing the current
linear OLS Scordilis (\citeyear{Scordilis2006}) coefficients with GOR-derived
relationships—either the published SHEEC coefficients of
\citet{Grunthal2016} or region-specific coefficients computed from the
overlap period of the source \catalogs—would reduce systematic bias in the
converted frequency-magnitude distribution without requiring any architectural
changes to the merging pipeline.

The second highest-priority gap is the adoption of a graph-based connected-
components step (see \cref{sec:graph}).  The current greedy forward-sweep misses
transitive associations and processes each event against only future events in
time order, meaning that when event B has already been consumed into a group
containing A, C cannot subsequently be merged with the A-B group even if C
matches B within threshold.  A Union-Find graph implementation would resolve
this in $O(n\,\alpha(n))$ time with no change to the downstream selection logic.
